% Complete checked source block A.5.9; source SHA256 b32705edb3aa3ed1f58c951bde13c7b6a0b1029b8c847158782902c19293fbe9
\begin{statement}{A.5.9}{Proved}{The ruled obstruction with a reduced conductor image}
Let $S\subset Y$ be an irreducible nonnormal reduced Cartier
divisor in a compact smooth complex threefold with
$H^2(Y,\Z)=0$ and $\chi(\OO_S)=0$.
Let $\eta:T\to S$ be its normalization and
$\pi:Z\to T$ its minimal resolution, and put $h=\eta\pi$.
Suppose that $Z$ is irrational ruled, with ruling $p:Z\to B$,
that every exceptional curve of $\pi$ is vertical, and that
the effective divisor of A.5.3 has the form
\[
 K_Z+D_Z=h^*\omega_S,\qquad D_Z=2C+V,                 \tag{A.RC.1}
\]
where $C$ is a section in the strict conductor and $V$ is
vertical. Suppose also that the conductor image scheme
$\Gamma=\operatorname{Spec}_S(\OO_S/\mathfrak c)$ is reduced,
and that $H^0(S,K\otimes L|_S)\ne0$ for every $L\in\Pic(Y)$,
with $K$ as in A.1.1.
Then every ambient line bundle pulls back trivially to $T$.
In particular
\[
 H^2\bigl(T,\omega_T\otimes
       \eta^*((K_Y^{-1}\otimes L)|_S)\bigr)\simeq\C
       \quad\hbox{for every }L\in\Pic(Y).            \tag{A.RC.2}
\]
The proof forces $B$ and $\Gamma$ to be smooth elliptic curves,
$Z=T$ to be minimal ruled, and $D_Z=2C$ with $C^2=0$.
Thus no local-complete-intersection assumption on the original
conductor image curve is needed in this reduced-scheme case.
\end{statement}
\begin{proof}
Write $\mathcal B=\eta_*\OO_T$ and
$\mathcal Q=\mathcal B/\OO_S$. We first specify conductor
duality, including its twists. Since $S$ is Gorenstein of
dimension two and $\mathcal B$ is maximal Cohen--Macaulay,
its dual is $\mathfrak c=\mathcal Hom_S(\mathcal B,\OO_S)$,
and maximal Cohen--Macaulay biduality gives
$\mathcal Hom_S(\mathfrak c,\OO_S)=\mathcal B$.
The quotient $\OO_\Gamma$ is Cohen--Macaulay of dimension one:
this follows from the depth lemma in
$0\to\mathfrak c\to\OO_S\to\OO_\Gamma\to0$.
Applying $\mathcal Hom_S(-,\OO_S)$ to that sequence, its
initial terms identify the actual normalization inclusion and
its cokernel as
\[
 0\longrightarrow\OO_S\longrightarrow\mathcal B
 \longrightarrow
 \mathcal Ext^1_S(\OO_\Gamma,\OO_S)\longrightarrow0,
 \qquad
 \mathcal Q\simeq\omega_\Gamma\otimes\omega_S^{-1}|_\Gamma.
                                                               \tag{A.RC.3}
\]
Adjunction makes $\omega_S$ the restriction of
$K_Y\otimes\OO_Y(S)$. Its degree is zero on every compact
curve, because $H^2(Y,\Z)=0$. Curve Riemann--Roch and Serre
duality therefore give
$\chi(\mathcal Q)=\chi(\omega_\Gamma)=-\chi(\OO_\Gamma)$.
On the other hand, A.5.8 gives rational singularities of $T$
and $\chi(\OO_T)=1-g(B)$. The normalization sequence and
$\chi(\OO_S)=0$ now give the exact numerical identity
\[
                   \chi(\OO_\Gamma)=g(B)-1.          \tag{A.RC.4}
\]

We next prove that $\Gamma$ is connected. Contract ruling
$(-1)$-curves from $Z$ to a minimal ruled bundle $Z_0\to B$.
Under a single contraction $\sigma:Z_1\to Z_2$ with exceptional
curve $E$, let $D_2=\sigma_*D_1$, let $a$ be the coefficient
of $E$ in $D_1$, and let $m$ be the multiplicity of $D_2$ at
the contraction point. The canonical divisor formula gives
\[
 K_{Z_1}+D_1=\sigma^*(K_{Z_2}+D_2)+(1+a-m)E.
                                                               \tag{A.RC.5}
\]
Numerical triviality and intersection with $E$ force $a=m-1$.
The same numerical triviality descends to $Z_2$. On $Z_0$
the pushed divisor is $2C_0+V_0$, with $C_0$ a section and
$V_0$ a nonnegative sum of fibres, so its support is connected.
Reverse the contractions. A centre outside the existing support
would produce coefficient $-1$ on its exceptional curve; that
coefficient persists on strict transforms and contradicts the
effectiveness of the final $D_Z$. A centre of multiplicity one
is a smooth point of one coefficient-one component and its
exceptional coefficient is zero, so it does not disconnect the
support. At a centre of multiplicity at least two the exceptional
curve has positive coefficient and connects all strict-transform
branches through that point. Induction proves that
$\operatorname{Supp}D_Z$ is connected.

Every positive exceptional component of $D_Z$ maps into
$\Gamma$. Indeed away from $\Gamma$ the surface $T$ equals
the Gorenstein surface $S$ and has rational singularities by
A.5.8. Such a singularity is crepant on its minimal resolution:
a generator of its dualizing sheaf pulls back holomorphically
by rational duality; its effective exceptional zero divisor
has nonnegative intersection with each exceptional curve by
adjunction and minimality, so negative definiteness makes it
zero, as in A.5.2. Thus (A.RC.1) has no positive exceptional
coefficient there. The strict conductor maps onto the whole
conductor support. It follows that
$h(\operatorname{Supp}D_Z)=\operatorname{Supp}\Gamma$,
which proves connectedness.

Let $D\subset\Gamma$ be the image of $C$. Localize (A.RC.3)
at the generic point of $D$. Because $\Gamma$ is reduced,
$\OO_\Gamma$ there is the function field $k(D)$, and its
dualizing module has dimension one over that field. The
elementary conductor quotient sequence is
\[
 0\longrightarrow\OO_\Gamma\longrightarrow
          \mathcal B/\mathfrak c
   \longrightarrow\mathcal Q\longrightarrow0.
\]
It follows that $\mathcal B/\mathfrak c$ has dimension two
over $k(D)$. The conductor coefficient two on $C$ contributes
$2\deg(C\to D)$ to this dimension. Therefore this degree is
one and there is no other normalization branch over $D$.
The generic transverse ring has conductor $t^2\C\{t\}$ and
is exactly $\C+t^2\C\{t\}=\C\{t^2,t^3\}$, the cusp ring.
Thus $D^\nu\simeq C\simeq B$. This deduction also rules out
any additional vertical branch over the same image curve.

A connected compact reduced curve has only constant global
holomorphic functions. From (A.RC.4) we therefore obtain
$h^1(\Gamma,\OO_\Gamma)=2-g(B)$.
Restriction to the reduced component $D$ gives a surjection
on $H^1$, because its kernel is a coherent sheaf on a curve
and has no $H^2$. Normalization of $D$ gives
$h^1(D,\OO_D)\ge g(B)$. Hence
\[
             2-g(B)\ge h^1(D,\OO_D)\ge g(B).         \tag{A.RC.6}
\]
As $g(B)\ge1$, all quantities equal one. The normalization
quotient of $D$ has length zero, so $D$ itself is smooth
elliptic. Statement C2d applies to this actual smooth curve
in $Y$ and makes the transversal cusp constant along it.
No other singular curve meets $D$. Connectedness of $\Gamma$
therefore gives $\Gamma=D$. Its normalization inverse image
is the single smooth curve $C$, the normalization is smooth
near it, and C2d gives $C^2=0$.

Adjunction in (A.RC.1) yields $V\cdot C=0$. A ruling
$(-1)$-curve outside $V$ cannot meet $C$, since its
intersection with $K_Z+2C+V$ would be at least $-1+2$.
A component of $V$ cannot meet $C$ either. Thus all
contractions to $Z_0$ avoid $C$ and preserve its square zero.
On the minimal bundle, $K_{Z_0}+2C_0$ is numerically trivial,
so its nonnegative vertical divisor $V_0$ is zero.
The first reverse blowup would be away from $C_0$ and have
coefficient $-1$ by (A.RC.5). There are no such blowups.
Thus $Z$ is minimal ruled and $D_Z=2C$.
An exceptional curve for $Z\to T$ would have negative square;
adjunction with $K_Z\equiv-2C$ would make it rational and
therefore vertical. A minimal ruled bundle has only fibres
of square zero, so no such exceptional curve exists and $Z=T$.

The constant-cusp computation A.2.4 now gives the injection
$K\hookrightarrow\eta_*(h^*\OO_Y(-S)\otimes\OO_Z(C))$.
Every ambient line is topologically trivial on $Z$ and is
the pullback of a degree-zero line on the elliptic curve $B$.
The section exact sequence, together with $R^1p_*\OO_Z=0$,
gives
\[
 0\longrightarrow\OO_B\longrightarrow p_*\OO_Z(C)
 \longrightarrow N_{C/Z}\longrightarrow0,
 \qquad\deg N_{C/Z}=0.                              \tag{A.RC.7}
\]
After a degree-zero twist $P$, both end terms have zero
sections unless $P=\OO_B$ or $P=N_{C/Z}^{-1}$.
The image of $\Pic(Y)$ in $\Pic^0(B)$ is connected:
$H^2(Y,\Z)=0$ makes $\Pic(Y)$ the connected image of
$H^1(Y,\OO_Y)$ under the exponential map.
All-twist persistence of $K$, and translation by the ambient
line $\OO_Y(-S)$, put this entire image in those at most two
points. It is therefore the identity. This proves triviality
of every ambient pullback. Serre duality on $T$ proves
(A.RC.2).
\end{proof}


% A.5.10 from part_A6_cusp_chern.tex



% Complete checked source block A.5.10; source SHA256 3e1170995ea0b81d414b33035e0ec19ecf2d787f953e0d2707ec34ca45ba96d2
\begin{statement}{A.5.10}{Proved}{Eliminating vertical conductor on a smooth ruled cusp normalization}
Let $Y$ be a smooth compact complex threefold with
$H^2(Y,\mathbb Q)=0$, let $S\subset Y$ be a reduced
irreducible divisor, and suppose its normalization
$h:Z\to S\subset Y$ is a smooth irrational
ruled surface $p:Z\to B$.
Suppose its conductor divisor is $D_Z=2C+V$, where
$C$ is a ruling section and $V$ is vertical, and
the generic transversal along $h(C)$ is an ordinary cusp.
Then $B$ is elliptic, $V=0$, $C^2=0$, and $Z$ is
already minimal.
The image $h(C)$ is smooth elliptic and $S$ is locally
an ordinary-cusp product along it. No local
complete-intersection hypothesis on the image curve is
required. The conductor-image scheme along the vertical
components may initially be nonreduced.
\end{statement}
\begin{proof}
Let $g=g(B)\geq1$, let $F$ be the full ruling fibre
class, and write $V=\sum_i c_i\Gamma_i$ for its
irreducible vertical components. Put
$v=V\cdot C$, $n=C^2$ and $k=2g-2$.
Conductor duality and Cartier adjunction give the exact
line-bundle comparison
\[
 K_Z\otimes\mathcal O_Z(2C+V)
       =h^*(K_Y\otimes\mathcal O_Y(S)).
                                                        \tag{A.J.21}
\]
Its right side has rational first Chern class zero.
Intersecting with $C$ and using adjunction on the section
gives
\[
 n=2-2g-v=-k-v.                                  \tag{A.J.22}
\]
Let $D_p$ be the divisorial common zero of
$dp:T_Z\to p^*T_B$. Its remaining zero scheme is
zero-dimensional, say $\Xi$, so the image is
$I_\Xi\otimes p^*T_B(-D_p)$.
The kernel $T_p$ is an invertible sheaf, by the depth
lemma. Taking determinants gives the actual comparison
\[
 T_p=K_Z^{-1}\otimes p^*K_B\otimes\mathcal O_Z(D_p).
\]
The section makes $dp$ surjective at every point of $C$.
Consequently $D_p$ and $\Xi$ are disjoint from $C$,
and $T_p|_C=N_{C/Z}$.

Put $\mathcal F=h^*T_Y$. Poincare duality on the compact
oriented six-manifold $Y$ also gives
$H^4(Y,\mathbb Q)=0$. Therefore
\[
 c_1(\mathcal F)=0,\qquad c_2(\mathcal F)=0
        \quad\hbox{in rational cohomology}.       \tag{A.J.23}
\]
The differential is generically injective. Let
$L\subset\mathcal F$ be the saturation of the image
of $T_p\xrightarrow{dh}\mathcal F$.
It is invertible: a saturated rank-one subsheaf of a
locally free sheaf on a smooth surface is reflexive.
The map $T_p\to L$ has no vertical divisorial zero.
Indeed $h$ is finite, so no ruling fibre is contracted;
in characteristic zero its derivative along that fibre
is nonzero generically.

At the generic point of $C$, an ordinary cusp has
normalization differential equivalent to that of
$(u,s)\mapsto(u^2,u^3,s)$.
The ruling tangent either lies in its kernel along $C$
or does not. The vanishing order in the first case is
exactly one, and in the second case it is zero.
Write $\epsilon=1$ or $\epsilon=0$ for these respective
alternatives. There are no other horizontal conductor
components, and the normalization is an immersion away
from the conductor. Hence the actual line comparison is
\[
 L=T_p\otimes\mathcal O_Z(\epsilon C).
                                                        \tag{A.J.24}
\]
This checks both positions of the derivative kernel
relative to the ruling tangent.

The quotient $E=\mathcal F/L$ is torsion-free of rank
two. The differential induces a generically nonzero map
\[
 T_Z/T_p=I_\Xi\otimes p^*T_B(-D_p)
                    \longrightarrow E\longrightarrow E^{**}.
\]
It extends uniquely to a map from $p^*T_B(-D_p)$ into
$E^{**}$, since both are locally free and the omitted
set has codimension two; locally this is Hartogs
extension of its coefficient functions.
Let $M$ be its saturated image line in the rank-two
bundle $E^{**}$. At the generic point of $C$, the
two-by-two differential minors have common vanishing
order one. Removing the order $\epsilon$ already
assigned to $L$ leaves order $1-\epsilon$ for this map.
At the generic point of $\Gamma_i$, take a plane transversal
to its image and a normalization parameter $u$ for the
corresponding plane-curve branch. If $m_i$ is its
multiplicity, its derivative modulo the nonzero tangent
along $\Gamma_i$ vanishes to order $r_i=m_i-1$.
These are all codimension-one zeros. Put
$R=\sum_i r_i\Gamma_i$ and $r=R\cdot C$. We obtain
\[
 M=p^*T_B(-D_p)\otimes
           \mathcal O_Z\bigl((1-\epsilon)C+R\bigr).
                                                        \tag{A.J.25}
\]

The plane-curve conductor gives the inequalities
\[
 2R\leq V,\qquad 2r\leq v.                      \tag{A.J.26}
\]
To spell out their local length content, on a single
plane branch of multiplicity $m_i$ the classes
$u,u^2,\ldots,u^{m_i-1}$ are linearly independent in
the normalization quotient. Its delta invariant is thus
at least $m_i-1$. The branch is Gorenstein. Finite
duality identifies the length of its ring modulo its
conductor with the length of its normalization quotient;
the intrinsic conductor exponent is consequently twice
that delta invariant. Other branches add their
nonnegative intersection multiplicities to this exponent.
For the last comparison, writing the plane equation
as $\prod_j f_j$ multiplies the residue denominator on
branch $i$ by $\prod_{j\ne i}f_j$, whose valuation is
exactly the sum of those intersection multiplicities.
This proves (A.J.26) for every branch, including cases
with nonreduced conductor-image scheme.

Let $\ell_1$ be the length of $E^{**}/E$. Saturation
of $M$ gives a zero-dimensional subscheme $\Xi_M$ and a
line bundle $N$ with
\[
 0\longrightarrow M\longrightarrow E^{**}
       \longrightarrow I_{\Xi_M}\otimes N\longrightarrow0.
\]
Write $\ell_2=\operatorname{length}\Xi_M$.
First Chern classes give $c_1(N)=-c_1(L)-c_1(M)$.
The second Chern classes in these typed quotient
sequences and (A.J.23) give the intersection identity
\[
 L^2+L\cdot M+M^2=\ell_1+\ell_2\geq0.            \tag{A.J.27}
\]
Indeed $c_2(E)=L^2$,
$c_2(E^{**})=L^2-\ell_1$, and
$c_2(E^{**})=-L\cdot M-M^2+\ell_2$.

Set $W=V+D_p$ and $B_v=R-D_p$, both vertical
divisors, the first effective. The exact comparisons
above and (A.J.21) give the numerical classes
\[
 L\equiv(2+\epsilon)C+W+kF,\qquad
 M\equiv(1-\epsilon)C+B_v-kF.
\]
The intersection form on vertical divisors is negative
semidefinite. One can check this directly by contracting
the ruling blowups: every vertical class is a multiple
of the pullback fibre class plus exceptional classes,
on which the form is negative definite. Hence
\[
 q_v=W^2+W\cdot B_v+B_v^2
      =(W+B_v/2)^2+3B_v^2/4\leq0.
\]
Its vanishing implies $W^2=B_v^2=0$; these classes
are then numerical multiples of the fibre class.

For $\epsilon=1$, using $D_p\cdot C=0$ in the
left side of (A.J.27) gives
\[
 \ell_1+\ell_2=-6k-3v+3r+q_v
                  \leq-6k-3v/2.
\]
Nonnegativity, with $k\geq0$ and $v\geq0$,
forces $g=1$, $v=0$, $r=0$ and $q_v=0$.

For $\epsilon=0$, the same intersection calculation reads
\[
 0\leq\ell_1+\ell_2
       =-6k-2v+4r+q_v\leq-6k.                   \tag{A.J.28}
\]
It forces $g=1$, $\ell_1=\ell_2=0$, $2r=v$,
and $q_v=0$.
Thus $E=E^{**}$ is locally free, and
$L\to\mathcal F$ is a subbundle inclusion.
Along $C$, the composite
$T_C\to\mathcal F|_C\to E|_C$ vanishes generically,
because $dh$ has rank one there and its nonzero ruling
tangent image is $L|_C$. It therefore vanishes
everywhere. The exact subbundle sequence gives an
actual nonzero holomorphic line-bundle map
\[
 T_C\longrightarrow L|_C=T_p|_C=N_{C/Z}.
                                                        \tag{A.J.29}
\]
It is nonzero since the finite map $C\to h(C)$ has
nonzero derivative generically. Its Hom line has degree
$n-(2-2g)=-v$ by (A.J.22). A nonzero section on a
compact curve has nonnegative degree, so $v=0$ in this
alternative too.

In both alternatives, $v=r=0$ and $q_v=0$.
Thus $W$ is numerically a fibre multiple and
$W\cdot C=V\cdot C+D_p\cdot C=0$, so it is
numerically zero. An effective nonzero divisor on the
projective ruled surface has positive intersection
with an ample class. It follows that $W=0$ and hence
$V=D_p=0$. Now $B_v=R$ is effective, has square zero
and has zero intersection with $C$, so the same
argument gives $R=0$. Equation (A.J.22) gives $C^2=0$.
Finally $K_Z+2C$ is numerically trivial. A ruling
$(-1)$-curve $E_0$ would satisfy
$0=(K_Z+2C)\cdot E_0=-1+2C\cdot E_0$, impossible
because $C\cdot E_0$ is an integer. There are therefore
no ruling blowups and $Z$ is minimal.

Statement A.5.7 now applies everywhere along
the conductor: the normalization is smooth and its
conductor is exactly $2C$. It makes the image smooth
and the surface locally the cusp product. The image is
elliptic because $C\to h(C)$ is an isomorphism.
\end{proof}


% A.5.11 from part_A.tex



% Complete checked source block A.5.11; source SHA256 bd570049e500cae2e523f7202eba509d0093358d9c840cf619fd6da107d4e4f7
\begin{statement}{A.5.11}{Proved}{The genus of a horizontal conductor image from a local torus group}
Let $f:Y\to C_0$ be a proper holomorphic map from a smooth
compact complex threefold to a smooth curve, with connected
fibres and complex two-torus general fibre. Let $S$ be an
irreducible reduced fibre support, $\eta:T\to S$ its
normalization, and $\pi:Z\to T$ a resolution which is
irrational ruled, with ruling $p:Z\to B$.
Suppose every exceptional curve of $\pi$ is vertical, and
the strict conductor has just one horizontal irreducible
component $C$, which is a section of $p$.
Let $D=\eta\pi(C)$ and let $\Sigma=D^\nu$ be its normalization.
Then $g(\Sigma)\le1$.
In particular, when $C\to D$ has degree one, $B$ is elliptic.
This applies to the horizontal-cusp case of $D_Z=2C+V$;
it allows arbitrary multiplicity of the full fibre and
arbitrary singularities or nonreduced conductor image on
the finitely many vertical curves.
\end{statement}
\begin{proof}
There is nothing to prove when $\Sigma$ is rational, so
assume $g(\Sigma)\ge1$.
Statement A.5.8 gives a holomorphic ruling
$\widetilde p:T\to B$ with $p=\widetilde p\pi$.
Identifying $C$ with $B$ by its section map, normalization
of its nonconstant map to $D$ gives a finite surjective
holomorphic map $\phi:B\to\Sigma$.
Put $\psi=\phi\widetilde p:T\to\Sigma$.

We describe precisely the finite identifications needed to
make this map descend to $S$. Every vertical conductor
component of $T$ lies over one of finitely many points of
$B$, so $\psi$ is constant on it. Its image under $\eta$
has rational normalization: a vertical component is the
image of a rational curve on $Z$, and a nonconstant map
from that curve to its image normalization is finite.
Such an image cannot equal $D$, because $g(\Sigma)\ge1$.
Consequently these vertical images meet $D$ at finitely
many points. Away from those points and from the finitely
many singular points of $D$, every normalization preimage
of a point of $D$ lies on $C$, and its $\psi$ value is the
unique corresponding point of $\Sigma$.
At a point outside $D$ with more than one normalization
preimage, all those preimages lie on vertical conductor
components. There are no other horizontal conductor
components by hypothesis. Thus there is a finite set
$A\subset\Sigma$ such that, for every $s\in S$, either
$\psi(\eta^{-1}(s))$ is a singleton or it is contained
in $A$. One can take $A$ to contain the constant values
on all vertical conductor components and all points of
$\Sigma$ above the singularities of $D$ and its
intersections with the vertical image curves.

Choose an embedded finite tree $J$ in the underlying
oriented surface $\Sigma$ containing $A$, and let
$q:\Sigma\to\Sigma/J$ collapse that tree to one point.
Such a tree is obtained by joining the finite set by
arcs and retaining a spanning tree after subdividing
their finitely many crossings. Triangulate the surface
with $J$ a subcomplex. A contractible subcomplex is a
cofibration; contracting it gives a homotopy equivalence
of this CW pair. Concretely the homotopy extension
property extends a contraction of $J$ and supplies a
homotopy inverse to the quotient map. In particular
\[
              q_*:\pi_1(\Sigma)\xrightarrow{\sim}
                         \pi_1(\Sigma/J).          \tag{A.Pi.1}
\]
The map $q\psi$ is constant on every fibre of the finite
proper surjection $\eta$. A proper continuous surjection
between these compact Hausdorff spaces is a quotient
map, so there is a unique continuous map
\[
 a:S\longrightarrow\Sigma/J,
 \qquad a\eta=q\psi.                              \tag{A.Pi.2}
\]
Let $\delta:\Sigma\to D\subset S$ be the curve normalization
map. On its dense smooth locus, lifting through $C$ gives
$a\delta=q$. Continuity gives this equality everywhere.
Equation (A.Pi.1) therefore makes $a_*$ surjective on
fundamental groups, without any assertion that a finite
point-identification quotient has the same fundamental
group as $\Sigma$.

We next extend the actual map (A.Pi.2) to a neighbourhood
of $S$ in $Y$. The simultaneous analytic triangulation
theorem of \cite[\S3, Theorem~2, p.~464]{Lojasiewicz1964}
triangulates the pair $(Y,S)$, since a complex analytic
subset is real semianalytic. Compactness makes the
subcomplex representing $S$ finite. After barycentric
subdivision it is full: vertices lying in it span a
face in that same subcomplex. Write $\lambda_v$ for
barycentric coordinates and let $L$ be its vertex set.
On the open neighbourhood where $\sum_{v\in L}\lambda_v>0$,
the formula
\[
 r(x)=\frac{\sum_{v\in L}\lambda_v(x)v}
                  {\sum_{v\in L}\lambda_v(x)}       \tag{A.Pi.3}
\]
defines a continuous retraction to $S$ in triangulation
coordinates. Its value lies in the face spanned by the
$L$ vertices of the same simplex, so these formulas
agree across faces. This supplies an actual open
neighbourhood $N$ and a map $r:N\to S$ fixed on $S$.

Let $c=f(S)$. Properness allows a small disc $\Delta$
about $c$ with $U=f^{-1}(\Delta)\subset N$ and no other
critical value. The map $(ar)|_U:U\to\Sigma/J$ is
surjective on fundamental groups: its restriction to
$\delta(\Sigma)\subset S\subset U$ is $q$, whose
fundamental-group map is an isomorphism.
Set $U^*=U\setminus S=f^{-1}(\Delta^*)$.
The inclusion induces a surjection
$\pi_1(U^*)\twoheadrightarrow\pi_1(U)$.
Indeed a loop in the smooth six-manifold $U$ can be
perturbed relative to its base point to be transverse
to an analytic stratification of $S$; every stratum
has real codimension at least two, so the resulting
one-dimensional loop misses every stratum.

Over $\Delta^*$, the proper submersion is a smooth
torus bundle. The fundamental-group exact sequence
and $\pi_2(\Delta^*)=0$ give
\[
 1\longrightarrow\Z^4\longrightarrow\pi_1(U^*)
       \longrightarrow\Z\longrightarrow1.
\]
A lift of the generator of $\Z$ splits this group
extension, hence
\[
 \pi_1(U^*)\simeq\Z^4\rtimes_T\Z
 \twoheadrightarrow\pi_1(U)
 \twoheadrightarrow\pi_1(\Sigma/J)
 \simeq\pi_1(\Sigma).                             \tag{A.Pi.4}
\]
The first group is metabelian: its commutator subgroup
is contained in the abelian normal subgroup $\Z^4$.
Every quotient in (A.Pi.4) is therefore metabelian.
A closed orientable surface of genus at least two
has a fundamental group surjecting onto the free
group on two generators: in its usual surface
presentation, retain two $a$ generators and send all
$b$ generators and the other $a$ generators to the
identity. That free group surjects onto $S_4$ by
sending its generators to $(1234)$ and $(12)$;
their conjugates supply the adjacent transpositions.
The commutator subgroup of $S_4$ is the nonabelian
group $A_4$, so this quotient is not metabelian.
This contradicts (A.Pi.4). Thus $g(\Sigma)=1$ in the
positive-genus case. When $C\to D$ has degree one,
$\Sigma\simeq C\simeq B$, which proves the final
assertion.
\end{proof}


% A.5.12 from part_A.tex



% Complete checked source block A.5.12; source SHA256 e5075f4e99742eaf0c520fbe9a62ab7d15f1833aa068cb6dd4665803f8016b4a
\begin{statement}{A.5.12}{Proved}{The arithmetic genus of conductor images from coherent duality}
Let $f:Y\to C_0$ be a proper connected-fibre map from a
smooth compact complex threefold with $H^2(Y,\Z)=0$ to
a smooth curve, and let its smooth general fibre have
trivial canonical bundle. This includes complex
two-tori and primary Kodaira surfaces.
Let $mS$ be a full fibre with irreducible
reduced nonnormal support $S$. Write $\eta:T\to S$ for
normalization and $\pi:Z\to T$ for a resolution.
If the pullback homomorphism
$\Pic(Y)\to\Pic(Z)$ is nonconstant, then every reduced
irreducible component $D$ of the conductor image scheme
has arithmetic genus at most one.
In particular every such component whose normalization
has positive genus is a smooth elliptic curve.
The conductor image scheme itself is allowed to be
nonreduced, and no local-complete-intersection assumption
is imposed on its reduced components.
\end{statement}
\begin{proof}
First we determine exactly the needed canonical line on
the reduced support, without a multiplicity formula for
the direct image. The sheaf $f_*K_Y$ is torsion-free of
generic rank one on the smooth curve and hence is a
line bundle. Choose a frame near $c=f(S)$ and evaluate
it in $K_Y$. A coordinate $t$ on the base disc supplies
the conormal frame $dt$, and adjunction identifies
$K_Y|_F$ with $K_F$ after dividing by this frame.
Under that line-bundle isomorphism, the evaluated
section restricts on a general smooth fibre to a
nonzero canonical two-form, which is nowhere zero.
Its zero divisor is therefore
vertical. Shrinking the base disc to remove the other
vertical zero divisors, irreducibility of the support
gives an effective divisor $kS$, with $k\ge0$, and an
actual line-bundle comparison on $U=f^{-1}(\Delta)$:
\[
 K_Y|_U\simeq\OO_U(kS),\qquad
 \omega_S\simeq\OO_Y((k+1)S)|_S,\qquad
 \omega_S^{\otimes m}\simeq\OO_S.                 \tag{A.CD.1}
\]
The last isomorphism uses the local base equation whose
zero divisor is exactly the full fibre $mS$.
In particular no assertion $k=m-1$ is used.

Hirzebruch--Riemann--Roch in
$0\to\OO_Y(-S)\to\OO_Y\to\OO_S\to0$ gives
$\chi(\OO_S)=0$: the rational Chern character of
$\OO_Y(-S)$ equals that of $\OO_Y$, since its first
Chern class is zero. A nonzero section of the torsion
line $\omega_S$ trivializes it. Indeed its $m$th power
is a nonzero holomorphic function on the compact
connected reduced space and hence is a nonzero constant.
Thus, with $e=h^0(S,\omega_S)$, we have
\[
 e\in\{0,1\},\qquad
 e=1\Longleftrightarrow\omega_S\simeq\OO_S,
 \qquad h^1(S,\OO_S)=1+e.                         \tag{A.CD.2}
\]
Here $h^0(\OO_S)=1$ and Serre duality identifies
$h^2(\OO_S)$ with $e$.

Let $\mathcal Q=\eta_*\OO_T/\OO_S$ and let
$\Gamma\subset S$ be the conductor image scheme.
Both $S$ and $T$ are connected and compact, so their
global holomorphic functions are the same constants.
The normalization sequence consequently gives an
injection and the indicated kernel:
\[
 H^0(S,\mathcal Q)\hookrightarrow H^1(S,\OO_S),
 \qquad
 \operatorname{im}H^0(S,\mathcal Q)
 =\ker\bigl(H^1(S,\OO_S)\to H^1(T,\OO_T)\bigr).
                                                               \tag{A.CD.3}
\]
Because $H^2(Y,\Z)=0$, $\Pic(Y)$ is the connected
exponential image of $H^1(Y,\OO_Y)$. A nonconstant
holomorphic group homomorphism out of it has nonzero
differential. The differential of the stated Picard
pullback is the actual composite
\[
 H^1(Y,\OO_Y)\longrightarrow H^1(S,\OO_S)
 \longrightarrow H^1(T,\OO_T)
 \longrightarrow H^1(Z,\OO_Z).
\]
Thus its middle arrow has positive rank. Equations
(A.CD.2)--(A.CD.3) imply $h^0(\mathcal Q)\le e$.

The conductor is Cohen--Macaulay of dimension one,
and the exact duality comparison (A.RC.3) applies
without a reducedness assumption on $\Gamma$:
$\mathcal Q=\omega_\Gamma\otimes\omega_S^{-1}|_\Gamma$.
Serre duality on that curve therefore gives
\[
 h^1(\Gamma,\omega_S|_\Gamma)
       =h^0(S,\mathcal Q)\le e.                   \tag{A.CD.4}
\]
For any reduced irreducible component $D$ of $\Gamma$,
restriction of the line gives an exact quotient
$\omega_S|_\Gamma\to\omega_S|_D$. Its kernel is a
coherent sheaf on a one-dimensional compact space and
has no $H^2$, so the corresponding map on $H^1$ is
surjective. Hence
$h^1(D,\omega_S|_D)\le e$.

For $e=1$, (A.CD.2) makes this precisely
$p_a(D)=h^1(D,\OO_D)\le1$.
For $e=0$, the same inequality makes
$h^1(D,\omega_S|_D)=0$.
The restriction $\omega_S|_D$ is torsion by (A.CD.1)
and has degree zero. Curve Riemann--Roch now gives
\[
 0\le h^0(D,\omega_S|_D)
       =\chi(D,\omega_S|_D)=1-p_a(D),
\]
again proving $p_a(D)\le1$.
Finally, for a reduced irreducible curve the
normalization exact sequence gives
$p_a(D)=g(D^\nu)+\operatorname{length}
(\nu_*\OO_{D^\nu}/\OO_D)$.
Positive normalization genus therefore forces genus
one and zero normalization quotient. A normal complex
curve is smooth, which proves the final assertion.
\end{proof}


% A.5.13 from part_A.tex



% Complete checked source block A.5.13; source SHA256 eef7475a6f1d8637fe522d47cd4c396f955afbd40bcb3992ac6a42e3fe81bc6a
\begin{statement}{A.5.13}{Proved}{The conductor bound when the general fibre has $q=1$ and $p_g=0$}
Let $f:Y\to\PP^1$ be a proper connected-fibre map from
a smooth compact complex threefold with
$H^1(Y,\Z)=H^2(Y,\Z)=0$. Suppose its smooth general
fibre has $h^1(\OO)=1$ and $h^2(\OO)=0$.
Then
\[
 R^1f_*\OO_Y=\OO_{\PP^1},\qquad R^2f_*\OO_Y=0,
 \qquad\Pic(Y)=\C\beta,                            \tag{A.Q1.1}
\]
where the last identification uses the exponential
coordinates of D.1.9 and its nonzero transgression
class $\beta$. Every full fibre has structure-sheaf
cohomology dimensions $(1,1,0)$.
For an irreducible reduced support $S$, every
nonconstant ambient Picard pullback to a normalization
resolution forces every reduced conductor component
to have arithmetic genus at most one. In particular
each positive-genus component is smooth elliptic,
without any lci or conductor-reducedness hypothesis.
\end{statement}
\begin{proof}
By E4, $E=R^1f_*\OO_Y$ is a line bundle. The section
induced by D.1.9's integral transgression class is
nowhere zero: its image under the natural base-change
map at each point is the nonzero coefficient class on
that full fibre. A zero of the section in $E$ would
have zero image under that map. Thus $E\simeq\OO$.
Coherent Leray now gives $H^1(Y,\OO_Y)=\C\beta$.
The exponential sequence and the two stated integral
cohomology vanishings give the Picard identification
in (A.Q1.1).

The coherent sheaf $G=R^2f_*\OO_Y$ has generic rank
zero, hence is a torsion sheaf. Riemann--Roch on $Y$
gives $\chi(\OO_Y)=\int_Yc_1c_2/24=0$, since
$c_1(Y)$ is zero in rational cohomology. The Euler
characteristic of coherent Leray is therefore
\[
 0=\chi(\OO_{\PP^1})-\chi(E)+\chi(G)
      =\operatorname{length}G.
\]
Thus $G=0$. Top-degree base change gives
$H^2(F,\OO_F)=0$ for every full Cartier fibre $F$.
The multiplication-by-parameter sequence then gives
$H^1(F,\OO_F)=E\otimes\C(c)$, of dimension one;
E4 gives $h^0(F,\OO_F)=1$.

For a full fibre $F=mS$ with irreducible reduced
support, the quotient $\OO_F\to\OO_S$ induces a
surjection on $H^2$, because its kernel has support
of dimension two and has no $H^3$.
Hence $h^2(S,\OO_S)=0$. Riemann--Roch in the Cartier
sequence on $Y$ gives $\chi(\OO_S)=0$ as in A.5.12,
so $h^1(S,\OO_S)=1$.

Put $\mathcal Q=\eta_*\OO_T/\OO_S$ for the
normalization $T$ and write $\Gamma$ for the conductor
image scheme. Nonconstant Picard pullback has nonzero
differential factoring through
$H^1(S,\OO_S)\to H^1(T,\OO_T)$.
The source has dimension one, so this map has zero
kernel. The normalization sequence (A.CD.3) gives
$H^0(\mathcal Q)=0$. Conductor duality and Serre
duality give
\[
 H^1(\Gamma,\omega_S|_\Gamma)=0,
 \qquad H^1(D,\omega_S|_D)=0                      \tag{A.Q1.2}
\]
for each reduced irreducible component $D$, where
the second map is the restriction surjection in the
proof of A.5.12. Adjunction makes $\omega_S$ ambient,
so $\deg(\omega_S|_D)=0$ by $H^2(Y,\Z)=0$.
Curve Riemann--Roch yields
$0\le h^0(D,\omega_S|_D)=1-p_a(D)$.
The normalization genus and its nonnegative delta
length then give the asserted smooth elliptic
conclusion whenever that genus is positive.
\end{proof}


% A.5.14 from part_A6_critical_curves.tex



% Complete checked source block A.5.16; source SHA256 a302573aec5d6c61d35d44973699bce8f65eef52302ccc8e567b7e5c995bd294
\begin{statement}{A.5.16}{Proved}{Rational normalization bounds with the exact constant-sheaf map}
Let $f:Y\to B$ be a proper holomorphic map from a smooth compact
complex threefold with $b_2(Y)=0$, with connected fibres and smooth
general fibre a two-torus. Let $mS$ be a full fibre with irreducible
nonnormal reduced support. Write $\eta:T\to S$ for normalization,
$\pi:Z\to T$ for its minimal resolution, and $\Gamma\subset S$
for the conductor-image scheme. Suppose that $Z$ is rational.
Put $Q=\eta_*\OO_T/\OO_S$ and $e=h^0(S,\omega_S)$.
Then $e\in\{0,1\}$, and
\[
\begin{gathered}
 R^1\pi_*\OO_Z=0,\qquad
 H^1(T;\mathbb C)=H^1(T,\OO_T)=H^2(T,\OO_T)=0,\\
 H^0(S,Q)\xrightarrow{\ \simeq\ }H^1(S,\OO_S),\qquad
 h^0(S,Q)=1+e,\quad h^1(S,Q)=e,\quad
 \chi(\Gamma,\OO_\Gamma)=-1.
\end{gathered}
\tag{A.RT.1}
\]
In particular $T$ has rational singularities. The coefficient map
\[
 H^1(S;\mathbb C)\longrightarrow H^1(S,\OO_S)
\tag{A.RT.2}
\]
is injective, so $b_1(S)\le1+e\le2$.

There is an effective integral divisor $D$ on $Z$ with
$\OO_Z(D)\simeq K_Z^{-1}$, whose strict-transform part is the
normalization conductor. It satisfies
\[
 h^0(D,\OO_D)=h^1(D,\OO_D)=1,
 \quad D_{\rm red}\text{ connected},
 \quad p_a(D_{\rm red})\le1.
\tag{A.RT.3}
\]
The conductor image is connected,
$p_a(\Gamma_{\rm red})\le2$, and
$e(\Gamma_{\rm red})\ge0$. Set-theoretically
$D_{\rm red}=\pi^{-1}(\eta^{-1}(\Gamma_{\rm red}))$.
Exceptional curves over points outside the conductor have
coefficient zero in $D$; they are not to be counted as conductor
branches. With $E_p=(\pi^{-1}(p))_{\rm red}$ the exact Euler
formula is
\[
 e(S)=e(Z)-e(D_{\rm red})+e(\Gamma_{\rm red})
       -\sum_{p\in\operatorname{Sing}T\setminus\eta^{-1}(\Gamma)}
                  \bigl(e(E_p)-1\bigr).
\tag{A.RT.4}
\]
In the torus smoothing of the statement, A.5.17 excludes
every point in this sum. Hence $T$ is smooth outside the
normalization conductor, and the sum is zero.
\end{statement}
\begin{proof}
The canonical evaluation argument in A.5.12 gives
$\omega_S^m\simeq\OO_S$ and $\chi(S,\OO_S)=0$.
Thus $e=0$ or $1$, with $e=1$ precisely when
$\omega_S\simeq\OO_S$, and
$h^1(S,\OO_S)=1+e$.

We first prove a useful fact without the rationality assumption.
Finite duality identifies the normal dualizing sheaf with
\[
 \omega_T=\mathfrak c_T\otimes\eta^*\omega_S
       \ \subset\ \eta^*\omega_S,
\tag{A.RT.5}
\]
where $\mathfrak c_T$ is the conductor ideal on $T$.
A nonzero global section of $\omega_T$ is therefore a nonzero
section of the torsion line $\eta^*\omega_S$ vanishing on its
nonempty conductor. Its $m$th tensor power would be a nonzero
holomorphic function on compact connected $T$ vanishing there,
which is impossible. Consequently
$H^2(T,\OO_T)=H^0(T,\omega_T)^\vee=0$.

For the proper resolution $\pi$, normality gives
$\pi_*\OO_Z=\OO_T$, and its fibres have dimension at most one.
The low-degree Leray sequence is
\[
\begin{split}
0\longrightarrow H^1(T,\OO_T)\longrightarrow H^1(Z,\OO_Z)
 \longrightarrow H^0(T,R^1\pi_*\OO_Z)\\
 \longrightarrow H^2(T,\OO_T)\longrightarrow H^2(Z,\OO_Z)
 \longrightarrow0.
\end{split}
\tag{A.RT.6}
\]
Rationality of $Z$ makes both its positive-degree groups zero.
Together with (A.RT.5), this gives $H^1(T,\OO_T)=0$ and
$R^1\pi_*\OO_Z=0$, since this last sheaf has finite support.
The fibres of $\pi$ are connected. Proper base change for the
constant sheaf gives $R^0\pi_*\mathbb C_Z=\mathbb C_T$.
The degree-one topological Leray edge map is consequently an
injection $H^1(T;\mathbb C)\hookrightarrow H^1(Z;\mathbb C)=0$.

In the normalization exact sequence
\[
 0\longrightarrow\OO_S\longrightarrow\eta_*\OO_T
       \longrightarrow Q\longrightarrow0,
\tag{A.RT.7}
\]
the map on $H^0$ is the identity on constants. The preceding
vanishings therefore make its connecting map
$H^0(Q)\to H^1(\OO_S)$ an isomorphism and its next connecting
map $H^1(Q)\to H^2(\OO_S)$ an isomorphism. These are the
maps and dimensions in (A.RT.1). In particular $\chi(Q)=1$.
The exact conductor-duality isomorphism proved in A.5.9 is
\[
 Q\simeq\omega_\Gamma\otimes\omega_S^{-1}|_\Gamma.
\tag{A.RT.8}
\]
The conductor image is a proper Cohen--Macaulay curve, including
any nilpotent structure. Serre duality on it gives
$\chi(Q)=-\chi(\omega_S|_\Gamma)$. A torsion line bundle has
degree zero on every reduced component, and therefore also has
total degree zero on this possibly nonreduced curve. Curve
Riemann--Roch gives
$\chi(\omega_S|_\Gamma)=\chi(\OO_\Gamma)$.
This proves $\chi(\OO_\Gamma)=-1$.

Here is the asserted injectivity of the complex coefficient map.
In the category of sheaves of complex vector spaces set
\[
 \mathcal V=\operatorname{coker}
     (\mathbb C_S\longrightarrow\eta_*\mathbb C_T).
\]
The diagram from constants to holomorphic functions induces the
specific sheaf map
\[
 \mathcal V\longrightarrow Q,
 \quad[(a_1,\ldots,a_r)]\longmapsto[(a_1,\ldots,a_r)] .
\tag{A.RT.9}
\]
At a point $s\in S$, the $r$ entries refer to the connected
normalization germs above $s$. A tuple of branchwise constants
that comes from an element of $\OO_{S,s}$ has equal values on
all these germs, because evaluation at their common image point
is one scalar. It is therefore a diagonal tuple. This proves
that (A.RT.9) is injective on every stalk. No $\OO_S$-module
structure on $\mathcal V$ is asserted or needed.

A finite map has zero positive-degree direct images of the
constant sheaf, by proper base change on its finite fibres.
Since $H^1(T;\mathbb C)=0$ and the map on constants is an
isomorphism, the connecting map from
$0\to\mathbb C_S\to\eta_*\mathbb C_T\to\mathcal V\to0$
identifies $H^0(S,\mathcal V)$ with $H^1(S;\mathbb C)$.
Naturality of the connecting maps gives the commutative square
\[
\begin{array}{ccc}
 H^0(S,\mathcal V)&\xrightarrow{\ \simeq\ }&H^1(S;\mathbb C)\\
 \big\downarrow&&\big\downarrow\\
 H^0(S,Q)&\xrightarrow{\ \simeq\ }&H^1(S,\OO_S).
\end{array}
\tag{A.RT.10}
\]
Its left arrow is injective by (A.RT.9), proving (A.RT.2).

A.5.3 gives the effective integral divisor $D$ and the exact
line-bundle relation
$K_Z\otimes\OO_Z(D)\simeq(\eta\pi)^*\omega_S$.
The torsion line on the right is trivial on the smooth rational
surface: its Picard group is torsion-free and has zero identity
component. Hence $\OO_Z(-D)\simeq K_Z$. The divisor sequence
$0\to K_Z\to\OO_Z\to\OO_D\to0$, rationality of $Z$ and
Serre duality prove (A.RT.3). In detail the first positive
cohomology of $K_Z$ is zero, and $H^2(K_Z)=\mathbb C$;
the positive cohomology of $\OO_Z$ is zero. Thus $D$ is
connected. Restriction $\OO_D\to\OO_{D_{\rm red}}$ is
surjective on $H^1$, because the kernel is supported in
dimension one. Consequently $p_a(D_{\rm red})\le1$.

We identify its support, rather than discarding exceptional
components. On a minimal resolution every exceptional
irreducible curve $E$ has $K_Z\cdot E\ge0$. Indeed
$E^2<0$, and adjunction gives
$K_Z\cdot E=2p_a(E)-2-E^2$; a smooth rational
$(-1)$-curve is excluded by minimality. If $E$ has coefficient
zero in $D$, then
$0=(K_Z+D)\cdot E$ forces $K_Z\cdot E=D\cdot E=0$.
Thus $E$ is disjoint from $D_{\rm red}$.
Every exceptional fibre over the normalization conductor meets
its strict transform. Connectedness of that fibre then forces
all its exceptional curves to have positive coefficient in $D$.
Conversely, a positive-coefficient exceptional curve over a point
outside the conductor would lie in a component of $D_{\rm red}$
disjoint from every strict conductor curve. The connectedness
just proved, and the nonempty strict conductor, exclude this.
This proves the stated support equality and the zero-coefficient
assertion. It also proves connectedness of $\Gamma_{\rm red}$.

For completeness the bound on its arithmetic genus follows
directly from (A.RT.8). One has
$h^1(\Gamma,\omega_S|_\Gamma)=h^0(Q)=1+e$,
and restriction to $\Gamma_{\rm red}$ is surjective on $H^1$.
For $e=1$ the line is trivial, giving
$p_a(\Gamma_{\rm red})\le2$. For $e=0$ its $H^1$ has
dimension at most one, and degree-zero Riemann--Roch gives
$p_a(\Gamma_{\rm red})=1-h^0(\omega_S|_{\Gamma_{\rm red}})
+h^1(\omega_S|_{\Gamma_{\rm red}})\le2$.
Finally normalization and resolution are isomorphisms off the
conductor and the finitely many exceptional fibres.
Stratifying by these disjoint closed sets gives (A.RT.4) by
additivity of the compactly supported Euler characteristic.
At a point outside the conductor, $T\to S$ is an
isomorphism and $S$ is normal. A normal complex surface
has isolated singularities. The reduced local equation
therefore has an isolated critical point whenever that
point is singular. Since the full fibre has equation
$u h^m$, the isolated clause of A.5.17 excludes it.
This proves the final assertion and eliminates the displayed
sum by an actual local comparison with the nearby torus.

We prove the asserted Euler bound for the conductor image,
including its possible positive-genus component. Put
$C=\Gamma_{\rm red}$. The normalization exact sequence for
the connected reduced curve $D_{\rm red}$ shows that the sum
of its normalized component genera is at most
$p_a(D_{\rm red})\le1$. Any component with genus one is
smooth, is unique, and meets the remaining rational curves
only in trees; a second attachment path joining distinct
points of that elliptic component would increase the
arithmetic genus to at least two. Each component of $C$ is
the finite image of a strict conductor component of
$D_{\rm red}$. Riemann--Hurwitz therefore shows that its
normalization is rational or elliptic, with at most one
elliptic-normalized component.

Write $n_0$ for the number of rational-normalized components
of $C$, $n_1\in\{0,1\}$ for the number of its
elliptic-normalized components, and $r_p,\delta_p$ for the
number of branches and the delta invariant at a singular
point. The exact coherent and topological normalization maps
give
\[
 \chi(\OO_C)=n_0-\sum_p\delta_p,
 \qquad e(C)=2n_0-\sum_p(r_p-1).
\]
For each reduced curve germ, independent constants on its
normalized branches give $\delta_p\ge r_p-1$.
The already proved arithmetic-genus bound and connectedness
give $\chi(\OO_C)\ge-1$, hence $e(C)\ge n_0-1$.
If $n_0\ge1$ this is nonnegative. If $n_0=0$, there is
one elliptic-normalized component and
$\sum\delta_p\le1$. A negative value could therefore only
be $-1$, from one ordinary node. Indeed its curve ring has
colength one in the direct sum of its two normalization
rings, and is contained in the subring of pairs with equal
constant value; equality of colengths makes these rings
equal. The two points over that node are distinct on the
elliptic normalization.

Here the hypersurface embedding excludes this last case.
At least one point $t\in T$ over the node receives both
local conductor branches. Otherwise partition the local
irreducible hypersurface branches according to which of
the two conductor branches they receive. Both parts are
nonempty. Two surface branches from opposite parts meet
in a curve, since their distinct equations in the regular
three-dimensional ambient local ring generate an ideal
of height two. That intersection curve belongs to the
conductor image, so is one of its two branches. Both
surface branches receive it, contradicting the partition.
For a locally irreducible hypersurface there is only one
normalization point and the assertion is immediate.

The connected resolution fibre $\pi^{-1}(t)$ lies in
$D_{\rm red}$ by the support equality proved above.
The strict conductor curves over the two branches are
elliptic, because a rational curve cannot map
nonconstantly to an elliptic curve. There is at most one
such curve on $D_{\rm red}$. If the two strict curves
were different, they would already give total genus two.
If they coincide, they meet $\pi^{-1}(t)$ at distinct
points: the map of that normal elliptic curve factors
through the normalization of its image, so one point
cannot map to both points over the node. The connected
resolution fibre then supplies an additional path joining
those points, outside the elliptic curve. Its components
are rational, since $T$ has rational singularities.
This produces a graph cycle in addition to the elliptic
genus, again forcing $p_a(D_{\rm red})\ge2$.
Both alternatives are impossible. Thus $e(C)\ge0$.
\end{proof}

We record the exact minimal-model calculation that follows from
A.5.16. This is a reduction of the outstanding Euler calculation;
the configurations with negative complement Euler number are
retained in the calculation.
Choose any sequence of contractions of smooth rational
$(-1)$-curves from $Z$ to a minimal rational surface $Z_0$,
and let $D_0$ be the pushforward of $D$. At one contraction
$\sigma:Z_i\to Z_{i-1}$, let $E$ be the exceptional curve,
$a$ its coefficient in $D_i$, and $r$ the multiplicity of
$D_{i-1}$ at the centre. The actual divisor comparison is
\[
 K_{Z_i}+D_i
   =\sigma^*(K_{Z_{i-1}}+D_{i-1})+(1+a-r)E.
 \tag{A.RT.11}
\]
Both canonical-plus-boundary line bundles are trivial, so
intersection with $E$ gives $a=r-1$. In particular the centre
lies on $D_{i-1}$, and effectivity is retained at every stage.
If $r=1$, the centre is a smooth point of a reduced boundary
component, $a=0$, and the strict boundary is isomorphic to the
old boundary. The blowup increases its complement Euler number
by one. If $r\ge2$, the new reduced boundary is exactly the
full inverse image of the old reduced boundary. Both the
surface and the boundary replace the centre by $\mathbf P^1$,
so the complement Euler number is unchanged. Thus
\[
 e(Z\setminus D_{\rm red})
 =e(Z_0\setminus (D_0)_{\rm red})+N_0,
 \qquad N_0\ge0,
 \tag{A.RT.12}
\]
where $N_0$ is the number of multiplicity-one centres in this
specified sequence of reverse blowups. Formula (A.RT.12) concerns
the resolution complement. For the actual torus smoothing,
the isolated clause of A.5.17 has already removed the
off-conductor exceptional correction in (A.RT.4). This local
argument is needed: on an abstract minimal $\mathbb F_2$,
its exceptional $(-2)$-section can be disjoint from an
anticanonical divisor without having arisen from one of the
$N_0$ blowups.

On $\mathbf P^2$, the divisor $D_0$ is an effective cubic.
Its possible reduced supports and their complement Euler
numbers are
\[
\begin{array}{c|c}
 \text{reduced support}&e(\mathbf P^2\setminus(D_0)_{\rm red})\\ \hline
 \text{smooth cubic, nodal cubic, cuspidal cubic}&3,\ 2,\ 1\\
 \text{line and conic, transverse or tangent}&1,\ 0\\
 \text{three lines forming a triangle, or concurrent}&0,\ -1\\
 \text{two distinct lines, or one line}&0,\ 1.
\end{array}
\tag{A.RT.13}
\]
For example the union of three concurrent lines has Euler
number $6-3+1=4$, whereas $e(\mathbf P^2)=3$.
The last row retains respectively a double line plus a line,
and a triple line. These are all cubic possibilities: an
irreducible singular cubic has arithmetic genus one and
therefore exactly a node or a cusp; the other degrees partition
three as $2+1$ or $1+1+1$.

On $\mathbb F_n$, with $n=0$ or $n\ge2$, write $C_0^2=-n$
for a minimal section and $F$ for a ruling fibre. Then
\[
 D_0\sim 2C_0+(n+2)F.
 \tag{A.RT.14}
\]
Its horizontal part has degree two over the ruling base.
If that part consists of two distinct sections $C_1+C_2$,
write $V$ for the vertical part, $d$ for its total degree and
$k=C_1\cdot C_2$. Adjunction on either section gives
$k+d=2$. Let $q$ be the number of distinct intersection
points of $C_1$ and $C_2$, and let $r$ count those points
whose ruling fibre occurs in $V_{\rm red}$. A vertical fibre
meets the two sections in two points unless they coincide at
that fibre. Additivity therefore gives the exact formula
\[
 e((D_0)_{\rm red})=4-q+r,
 \qquad e(\mathbb F_n\setminus(D_0)_{\rm red})=q-r\ge0.
 \tag{A.RT.15}
\]
This includes two sections and a fibre through their common
point: that union has Euler number four and complement Euler
number zero.

If the horizontal part is an irreducible bisection $C$, the
same adjunction calculation gives $p_a(C)=1-d$, so $d=0$
or $1$. For $d=0$, its Euler number is at most two, and the
complement Euler number is at least two. For $d=1$, the
bisection is a smooth rational curve, the vertical part is one
reduced fibre, and their set-theoretic intersection has one or
two points. The complement Euler number is respectively one
or two.

The remaining horizontal part is $2C$, where $C$ is a section.
Write $C\sim C_0+kF$. The vertical divisor has total degree
$d=n+2-2k$, and, with $v$ its number of distinct fibres,
\[
 e((D_0)_{\rm red})=2+v,
 \qquad e(\mathbb F_n\setminus(D_0)_{\rm red})=2-v.
 \tag{A.RT.16}
\]
If $C\ne C_0$, positivity of intersection with $C_0$ gives
$k\ge n$; hence $d\le2-n$ and $v\le2$.
For $n=0$ one also has $d\le2$. Consequently the only
negative values on a minimal ruled surface occur in the
explicit configurations
\[
 Z_0=\mathbb F_n,\quad n\ge2,\qquad
 D_0=2C_0+\sum_{j=1}^{v}a_jF_j,\qquad
 a_j\ge1,\quad \sum_j a_j=n+2,\quad v\ge3.
 \tag{A.RT.17}
\]
The fibres $F_j$ here are distinct. Equations
(A.RT.11)--(A.RT.17), together with (A.RT.4), isolate the
remaining contributions exactly; an arbitrary identification
of their conductor branches has not been substituted for an
actual hypersurface normalization map.

There is a further fundamental-group restriction on precisely
these negative minimal-model configurations. Their reduced
supports are unions of rational curves with simply connected
underlying spaces: three concurrent lines on $\mathbf P^2$,
or the section with its attached ruling fibres in (A.RT.17).
Every reverse blowup in (A.RT.11) preserves this property.
A multiplicity-one blowup leaves the support isomorphic to
the preceding one. A positive-coefficient blowup at a smooth
point adds a rational leaf, at an ordinary double point
subdivides the incidence edge, and at the initial concurrent
point replaces that vertex by an exceptional $\mathbf P^1$
with three attached branches. Subsequent supports have only
the resulting ordinary double points. None of these explicit
operations introduces a topological loop. Consequently
$D_{\rm red}$ is simply connected throughout these cases.

The actual map $h=\eta\pi$ expresses $S$ as the topological
pushout of
\[
 D_{\rm red}\hookrightarrow Z,
 \qquad D_{\rm red}\xrightarrow{\ h|_{D_{\rm red}}\ }
                           \Gamma_{\rm red}.
 \tag{A.RT.18}
\]
Indeed $h$ is an isomorphism off these closed sets and the
whole preimage of $\Gamma_{\rm red}$ is $D_{\rm red}$.
The continuous map from this pushout to $S$ is a bijection
from a compact space to a Hausdorff space, and hence is a
homeomorphism. A triangulation of the pair
$(Z,D_{\rm red})$ makes its inclusion a cofibration.
The double mapping-cylinder model for (A.RT.18) therefore
has the same homotopy type as the pushout, and van Kampen
gives
\[
 \pi_1(S)\simeq
 \pi_1(Z)*_{\pi_1(D_{\rm red})}\pi_1(\Gamma_{\rm red})
 \simeq\pi_1(\Gamma_{\rm red}).
 \tag{A.RT.19}
\]
All normalizations of the image components are $\mathbf P^1$,
by Riemann--Hurwitz. Gluing their finitely many points over
singularities identifies the topological curve with a graph
of simply connected spheres. Its fundamental group is a
free group, of rank $b_1(\Gamma_{\rm red})$.

That free group is a quotient of a metabelian group. To see
the exact comparison, triangulate $(Y,S)$ and retract a small
regular neighbourhood $N$ onto $S$, as in A.5.11. Properness
puts a sufficiently small full fibre tube $U=f^{-1}(\Delta)$
inside $N$. Its retraction is surjective on fundamental
groups because it restricts to the identity on $S$.
General position makes
$\pi_1(U\setminus S)\to\pi_1(U)$ surjective, and the
torus bundle over $\Delta^*$ has
$\pi_1(U\setminus S)=\mathbb Z^4\rtimes\mathbb Z$.
This last group has abelian commutator subgroup, so its
quotient (A.RT.19) is metabelian. A free group of rank at
least two is not metabelian: it maps onto $S_4$, whose
commutator group $A_4$ is not abelian. We obtain
\[
 \pi_1(S)\simeq\pi_1(\Gamma_{\rm red})\in\{1,\mathbb Z\},
 \qquad b_1(\Gamma_{\rm red})\le1.
 \tag{A.RT.20}
\]
Under the full CDP Betti hypotheses, the nonzero class on
every support from D.1.9 excludes the trivial group in
(A.RT.20). This last use of the global transgression is
separate from the preceding $b_2(Y)=0$ calculations.

For these two boundary classes one can further exclude a
one-component image with first Betti number one. Such an image
has one two-branch singular point $q$; write $\alpha,\beta$
for its two normalization points there. Any two normalization
points of $S$ above $q$ receive a common local conductor branch,
because the corresponding two hypersurface branches meet in
a curve. Thus there cannot be both a point receiving only
$\alpha$ and a point receiving only $\beta$. After interchanging
the labels, every normalization point above $q$ receives
$\alpha$.

Form the incidence graph with one vertex for each strict
conductor component and one vertex for each connected
resolution fibre over a normalization point above $q$.
Edges record their intersection points, labelled by the
branch $\alpha$ or $\beta$ received there. Smooth resolution
fibres, consisting of one point, are included. The graph is
a forest: it is obtained from the tree of components and
intersection points of $D_{\rm red}$ by contracting disjoint
exceptional subtrees and deleting the other incidence edges.
Every strict conductor component maps finitely onto the sole
image component, so it has a preimage of $\beta$ and therefore
a $\beta$-edge. Every resolution-fibre vertex has an
$\alpha$-edge by the preceding paragraph. Starting at any
strict-component vertex, follow a $\beta$-edge and then an
$\alpha$-edge, and repeat. This walk never immediately reverses
an edge, because its consecutive labels differ. It is an
infinite nonbacktracking walk in a finite forest, which is
impossible. The contradiction proves the exclusion.

Since every image component is rational-normalized, write
$n_\Gamma$ for their number and $\beta_\Gamma\le1$ for the
first Betti number. Their exact topological normalization
formula now yields, on these negative boundary classes,
\[
 e(\Gamma_{\rm red})=n_\Gamma+1-\beta_\Gamma\ge2.
 \tag{A.RT.21}
\]
For one component the excluded case is the only value below
two; for at least two components the inequality follows
directly from $\beta_\Gamma\le1$.


% A.5.17 from part_A6_milnor.tex



% Complete checked source block A.5.18; source SHA256 b15a24ff1d58a286d50534f39550909badda17e7322d208c12e30fd29e14f79b
\begin{statement}{A.5.18}{Proved}{Every horizontal tacnode image is elliptic}
Let $f:Y\to C_0$ be a proper holomorphic map from a smooth
compact complex threefold with $b_2(Y)=0$, with connected
fibres and smooth general fibre a complex two-torus.
Let $mS$ be a full fibre with irreducible reduced support,
$\eta:T\to S$ its normalization, and $\pi:Z\to T$ a
resolution with an irrational ruling $p:Z\to B$.
Suppose every exceptional curve of $\pi$ is vertical and
the conductor divisor on $Z$ is $2C+V$, where $C$ is a
section of $p$ and $V$ is vertical. Let $D=\eta\pi(C)$.
Suppose the generic plane-curve transversal to $D$ in $S$ is an
ordinary tacnode.

Then $D$ is a smooth elliptic curve, $B=C$ is elliptic,
and $C\to D$ is an unramified double cover. The transversal
type along $D$ is everywhere the same ordinary tacnode.
The normalization $T$ is a smooth minimal ruled surface
over this elliptic curve. Its conductor is twice a
section, and $e(S)=0$. In particular its minimal
resolution is the identity.
In particular a rational image of
the horizontal tacnodal conductor cannot occur. The proof
allows vertical conductor and rational singularities of
$T$; it does not assume that the normalization map is
immersive at the special points under examination.
\end{statement}
\begin{proof}
If the normalization of $D$ has positive genus, A.5.14
and A.5.5 make $D$ smooth elliptic and its transversal
type constant. There are two generic normalization
branches of a tacnode. A vertical component of $Z$
is rational, so it cannot dominate this elliptic curve.
Thus both branches lie on $C$, and $C\to D$ is an
unramified double cover: the constant local equation
$x^2-y^4=0$ has two separately smooth normalization
branches on every sufficiently small disc in $D$.
Riemann--Hurwitz gives $g(C)=g(B)=1$.
It remains to exclude the case $D^\nu=\Sigma\simeq\PP^1$.

The finite map $C\to\Sigma$ has degree at least two,
because $C\simeq B$ has positive genus. Since a generic
tacnode has exactly two normalization branches, this
degree is exactly two. Consequently no vertical
conductor component can also dominate $D$: its generic
points would supply an additional normalization branch.
Write $g=g(C)\ge1$ and let $\mathcal B\subset\Sigma$
be the reduced branch divisor of this double cover.
Riemann--Hurwitz gives
\[
                   \deg\mathcal B=2g+2.             \tag{A.T.1}
\]

We first bound the number of branches of the singular
image curve itself. A.5.8 supplies the descended ruling
$\widetilde p:T\to B$. Identify $C$ with $B$ using the
section, let $\phi:B\to D$ be its actual map to the
image curve, and put
\[
                    \psi=\phi\widetilde p:T\to D.
\]
Every vertical conductor component has constant $\psi$
value. Its image under $\eta$ differs from $D$, as just
proved, and meets $D$ in finitely many points. Thus
there is a finite set $A\subset D$ such that, for each
$s\in S$, either $\psi(\eta^{-1}(s))$ is a singleton
or all its values belong to $A$. Include in $A$ the
constant values on the vertical conductor components
and the values of $\phi$ over their finitely many
intersections with $D$. Fibres over the remaining
points of $D$ have their preimages on $C$, where all
the $\psi$ values equal that same point of $D$.

Choose an embedded finite tree $J\subset D$ containing
$A$, in a triangulation of the compact analytic curve.
It can be constructed by joining these finitely many
points by arcs, subdividing their intersections and
retaining a spanning tree. The quotient
$q:D\to D/J$ is a homotopy equivalence: a contractible
CW subcomplex is a cofibration, and the homotopy
extension property extends its contraction. Since
$q\psi$ is constant on every fibre of the finite
proper quotient $\eta$, there is a continuous map
\[
 a:S\to D/J,\qquad a\eta=q\psi,\qquad a|_D=q.     \tag{A.T.2}
\]
The last equality holds first on the dense regular
locus of $D$ and then everywhere by continuity.

Use the explicit triangulated neighbourhood retraction
$r:N\to S$ of (A.Pi.3). For a sufficiently small disc
$\Delta$ around the fibre value, properness gives
$U=f^{-1}(\Delta)\subset N$. The map $(ar)|_U$ is
surjective on fundamental groups because its restriction
to $D$ is the homotopy equivalence $q$. Perturbing loops
off the real-codimension-two analytic support gives
$\pi_1(U\setminus S)\twoheadrightarrow\pi_1(U)$.
The torus-bundle exact sequence over $\Delta^*$ therefore
gives the actual quotient maps
\[
 \Z^4\rtimes\Z\twoheadrightarrow\pi_1(U)
       \twoheadrightarrow\pi_1(D/J)\simeq\pi_1(D).
                                                        \tag{A.T.3}
\]
The first group is metabelian, so the last is metabelian.

For each singular point $x$ of $D$, let $r_x$ be the
number of its normalization branches and put
$\beta=\sum_x(r_x-1)$. The finite normalization map
identifies exactly those finite sets of points on
$\PP^1$; since it is a compact Hausdorff quotient,
this is also the actual underlying topological quotient.
Adjoining arcs between the points to be identified
gives a CW presentation homotopy equivalent to
$S^2\vee\bigvee^\beta S^1$. Hence
$\pi_1(D)$ is the free group $F_\beta$.
For $\beta\ge2$ this group surjects onto $S_4$ by
sending two free generators to $(1234)$ and $(12)$;
the commutator $A_4$ of that quotient is nonabelian.
Equation (A.T.3) therefore forces $\beta\le1$.
By A.5.14 the total delta invariant $\delta=p_a(D)$
is at most two. If $s$ denotes the number of points
of $\Sigma$ over $\operatorname{Sing}D$, then
\[
 s=\#\operatorname{Sing}D+\beta\le\delta+\beta\le3.
                                                        \tag{A.T.4}
\]

First suppose that $D$ is a local complete intersection.
We retain the exact bundles and jets of A.5.5.
Let $\nu:\Sigma\to D$ be normalization, $I=\mathcal I_D$,
\[
 N=(\nu^*(I/I^2))^*,\qquad
 M=\nu^*(\mathcal O_Y(S)|_D).
\]
Thus $M$ is a degree-zero line bundle on $\Sigma$.
Let $\Xi$ be the curve-normalization
conductor divisor, $c=\deg\Xi=2\delta$, and let
$R\subset N$ and $Q=N/R$ be the saturated radical and
quotient of the generic rank-one Hessian. Write
$q_0=\deg Q$. The adjunction and conductor comparisons
(A.J.1) give
\[
 \deg N=-2+c,\quad \deg R=-2+c-q_0,\quad
 0\ne a\in H^0(\Sigma,M\otimes Q^{-2}),\quad q_0\le0.
                                                        \tag{A.T.5}
\]
Here $a$ denotes the intrinsic quadratic coefficient,
whose zero divisor has degree $-2q_0$.
The actual quartic invariant $\Delta=ae-b^2/4$ of
A.5.5 is meromorphic in
$L^2$, where $L=M\otimes Q^{-1}\otimes R^{-2}$.
Its pole bound (A.J.9) gives a nonzero holomorphic
section $\sigma$ of the actual square line bundle
$L^2(4\Xi)=(L(2\Xi))^2$. Its zero divisor has degree
\[
               d_\Delta=8+2q_0.                    \tag{A.T.6}
\]

The double cover $C\to\Sigma$ is obtained by adjoining
$\sqrt{-\Delta}$ to the function field, with the line
trivializations specified by $L$. Indeed over the
generic point a tacnode has normal expansion
$a\ell^2+b\ell r^2+er^4$, and the two normalization
branches are distinguished by the two roots of
$a(\ell/r^2)^2+b(\ell/r^2)+e$. Their discriminant
is $-4\Delta$. Changing a local frame of $L$ multiplies
the representing function of $\Delta$ by a square,
so this describes the same degree-two extension on
every overlap. At a point with local parameter $z$,
write that representing function as $z^n u(z)$,
$u(0)\ne0$. A root of the unit and removal of the
even power of $z$ show that its normalization is
ramified exactly when $n$ is odd. Adding $4\Xi$ changes
each order by an even number. Thus $\mathcal B$ is
exactly the odd-multiplicity support of
$\operatorname{div}(\sigma)$, and
\[
                  2g+2\le 8+2q_0.                 \tag{A.T.7}
\]

At a point of $D_{\rm reg}$ where $a\ne0$, A.5.17
applies to the actual local $m$th root of the full-fibre
map. Its nonzero Hessian makes the local equation a
constant transverse product, necessarily of type
$A_3$ because this is the generic type along $D$.
The conductor cover is therefore unramified there.
Every branch point outside $\nu^{-1}(\operatorname{Sing}D)$
is consequently a zero of $a$. Equations (A.T.4)--(A.T.5)
give the second bound
\[
                    2g+2\le -2q_0+s.              \tag{A.T.8}
\]
Adding (A.T.7) and (A.T.8) yields
$4g+4\le8+s\le11$. Since $g\ge1$, we have
\[
           g=1,\qquad \deg\mathcal B=4,
                    q_0\in\{-2,-1,0\}.            \tag{A.T.9}
\]

We next prove the exact local obstruction required in
the equality cases. Let a smooth point of $D$ be a
branch point at which $\Delta$ has a simple zero.
The preceding nonzero-Hessian argument gives $a(0)=0$.
In coordinates $D=(x,y)=0$ with curve parameter $z$,
choose the holomorphic saturated Hessian factor as the
first normal coordinate. The quadratic jet is then
$a(z)x^2$. The radical cubic vanishes identically,
so a local reduced equation has the exact expansion
\[
 f_0=a(z)x^2+x^2(xA+yB)+xy^2b(z)+xy^3C+y^4E,
                                                        \tag{A.T.10}
\]
with holomorphic coefficients. The identity
$\Delta=ae-b^2/4$, $a(0)=0$ and the simple zero of
$\Delta$ imply $b(0)=0$, $a'(0)\ne0$ and $e(0)\ne0$.
Indeed an order at least two for $a$, together with
$b(0)=0$, would give an order at least two for $\Delta$;
and then $\Delta'(0)=a'(0)e(0)$ is nonzero.
Thus $b/a$ is holomorphic and $E(0)\ne0$.

Set $X=x+b(z)y^2/(2a(z))$. In this actual analytic
coordinate change the coefficient of $Xy^2$ cancels.
All terms created from $x^2(xA+yB)$ lie in
$(X^2)+(Xy^3)+(y^4)$. Therefore, collecting the
terms of $X$-degree zero and one, gives exactly
\[
 f_0=X^2A_1(X,y,z)+Xy^3B_1(y,z)+y^4C_1(y,z),
 \qquad C_1(0)\ne0.
\]
Choose a fourth root of $C_1$ and define
\[
 Y=yC_1^{1/4}+
                  \frac{XB_1}{4C_1^{3/4}}.
                                                        \tag{A.T.11}
\]
The map $(X,y,z)\mapsto(X,Y,z)$ is locally
biholomorphic, and direct expansion gives
$f_0-Y^4\in(X^2)$. Hence $f_0=Y^4+X^2H(X,Y,z)$.
Its quadratic jet along $D$ is still $a(z)X^2$,
so $H(0,0,z)=a(z)$ and $\partial H/\partial z(0)\ne0$.
The further coordinate $Z=H$ gives exactly
\[
                         f_0=Y^4+X^2Z.             \tag{A.T.12}
\]
To apply this to the actual map, first take a local
$m$th root of its unit. This multiplies the reduced
equation by a unit, which changes neither zero order
of $a$ nor of $\Delta$. The same coordinate construction
then makes the full map $(Y^4+X^2Z)^m$.
Statement A.5.15 excludes it. Consequently every
smooth branch point has order at least three in
$\operatorname{div}(\sigma)$.

We also record precisely when the conductor denominator
can be removed at a singular point $x\in D$.
If the intrinsic section $a$ is nonzero at every
normalization point above $x$, choose an actual
constant linear combination of the two lci conormal
generators which is nonzero under the normalized
quotient at all those finitely many points. The
corresponding diagonal entry $a_0$ of the symmetric
Hessian matrix belongs to $\mathcal O_{D,x}$ and
has nonzero value on every normalization branch.
It is therefore a unit in $\mathcal O_{D,x}$.
The other normalized factor coefficient is the
actual quotient $h=q_{XY}/a_0\in\mathcal O_{D,x}$.
Thus square completion and monic polynomial division
in the proof of (A.J.7)--(A.J.8) both take place
over $\mathcal O_{D,x}$ itself. The regular-sequence
graded pieces inject into their normalization
pullbacks, so their cancellations already descend
on $D$. We may repeat that proof with $d=1$.
It gives a holomorphic section $\Delta$ of $L^2$
above $x$, without a conductor pole. At a branch
point above $x$, this proves the sharper inequality
\[
 \operatorname{ord}\sigma\ge
                        4\operatorname{ord}\Xi+1.
                                                        \tag{A.T.13}
\]

We now exhaust the three integer cases in (A.T.9).
If $q_0=-2$, the divisor of $\sigma$ has degree four
and has four odd-multiplicity points, so each is
simple. At most three lie above singularities of
$D$, by (A.T.4). The remaining smooth branch point
contradicts (A.T.12).
If $q_0=0$, the section $a$ has no zeros and all
four branch points would have to lie above the
singular locus; this contradicts $s\le3$.

Finally let $q_0=-1$, so that $\deg\operatorname{div}(a)=2$
and $d_\Delta=6$. If $k$ of the four branch points
lie above $\operatorname{Sing}D$, then the local
obstruction and parity give
\[
               6\ge k+3(4-k)=12-2k.
\]
Hence $k=s=3$, and equality holds: the three singular
preimages have order one in $\operatorname{div}(\sigma)$
and the one smooth branch point has order three.
The relations $s\le\delta+\beta\le3$ now force
$\delta=2$, $\beta=1$ and two singular points of
delta invariant one. The exact delta-one
classification in A.5.14 identifies them as one
ordinary node and one cusp. Their normalization
conductor orders are respectively $(1,1)$ and $2$;
these follow directly from the displayed node and
cusp rings in that proof.

At the cusp preimage, if $a$ were nonzero,
(A.T.13) would make the order of $\sigma$ at least
$4\cdot2+1=9$, exceeding its total degree six.
So $a$ vanishes there. It also vanishes at the
smooth branch point. These two distinct points
exhaust the degree-two divisor of $a$. The section
is consequently nonzero at both node preimages.
Apply (A.T.13) at the node: each of its two branch
points has order at least $4\cdot1+1=5$ in
$\operatorname{div}(\sigma)$, again exceeding the
total degree six. This excludes every rational lci image.

We finish by treating all non-lci image germs permitted
by $\delta\le2$ and $\beta\le1$, with explicit analytic
rings and transverse coordinate maps. For a unibranch
reduced curve ring $A\subset\C\{t\}$ of colength at
most two, an element of order one makes $A=\C\{t\}$.
Otherwise its least positive valuation is two or three:
a least valuation at least four leaves three linearly
independent missing coefficients. If it is three and
the colength is two, the inclusion
$A\subset\C+t^3\C\{t\}$ is equality, since the latter
already has colength two. If it is two, choose a
parameter with $t^2\in A$ by taking a square root of
the leading unit. The even part of every convergent
series belongs to $\C\{t^2\}\subset A$. Subtracting
it shows that the least odd element can be written
$t^{2k+1}v(t^2)$; division by its unit puts
$t^{2k+1}$ in $A$. Consequently
$A=\C\{t^2\}\oplus t^{2k+1}\C\{t^2\}$, of colength
$k\le2$. These are the planar cusp and the planar
$A_4$ curve, besides the smooth case.

For two branches, write their image rings as $A_1,A_2$.
The two projections of $A\subset A_1\oplus A_2$ are
surjective. Their kernels define a common finite
quotient $Q_0$ and an exact fibre-product description
$A=A_1\times_{Q_0}A_2$. Thus
$\delta=\delta_1+\delta_2+\operatorname{length}Q_0$.
If both branches are smooth, the quotient is
$\C\{t\}/(t^k)$ with $k\le2$; after changing the
second parameter, the two maps identify the same
$k$-jet. The map
\[
 x\longmapsto(t,s),\qquad
 y\longmapsto(t^k,-s^k)
\]
identifies this fibre product with
$\C\{x,y\}/(y^2-x^{2k})$: write a matching pair as
the sum of its average and $y$ times its difference
divided by $2x^k$. These are the node and tacnode.
If one branch is singular, the delta budget forces
one cusp, one smooth branch, and $Q_0=\C$.
Three branches are excluded by $\beta\le1$.

An abstract plane curve ring in the given smooth
threefold is lci by the actual transverse-coordinate
construction in A.5.14. Thus the only remaining
non-lci analytic rings, necessarily consuming the
whole delta budget at a single point, are
\[
\begin{array}{ll}
 A=\C\{t^3,t^4,t^5\},
 &I=(y^2-xz,\ yz-x^3,\ z^2-x^2y),\\[2pt]
 A=\C\{t^2,t^3\}\times_{\C}\C\{s\},
 &I=(xz,\ yz,\ y^2-x^3).
\end{array}                                             \tag{A.T.14}
\]
These are normal forms of the actual embeddings.
Their three displayed coordinate functions form a
basis of the cotangent space of each ring. Lift them
to $\mathcal O_{Y,x}$; their differentials are linearly
independent because the embedding dimension is three.
They are therefore ambient coordinates by the inverse
function theorem. The displayed ideals are their
kernels, obtained respectively by the parametrizations
$(t^3,t^4,t^5)$ and the union of $(t^2,t^3,0)$ with
$(0,0,s)$.

We show that the local reduced equation $f_0$ actually
lies in $I^2$ in each case. In the first case put
\[
 g_1=y^2-xz,\quad g_2=yz-x^3,\quad
 g_3=z^2-x^2y,\quad
 F=x^5+xy^3+z^3-3x^2yz.
\]
The exact identities
\[
 xF=g_2^2-g_1g_3,\qquad
 yF=xg_1^2+g_2g_3,\qquad
 zF=g_3^2+xg_1g_2                              \tag{A.T.15}
\]
show that the maximal ideal annihilates the class of
$F$ modulo $I^2$. Moreover
$I^{(2)}=I^2+(F)$ in the convergent local ring.
Here is a finite-module verification of that equality.
Let $J=I^2+(F)$. The quotient $\C\{x,y,z\}/J$ is
finite over the discrete valuation ring $\C\{x\}$,
since its reduction is the finite parametrized curve.
Its generic rank is nine: the curve has degree three
over the $x$ coordinate, and the square of its smooth
codimension-two ideal has length three along each
generic normal slice. Modulo $x$, the quotient has
the nine-element basis
\[
 1,y,y^2,y^3,\ z,yz,y^2z,\ z^2,yz^2,
\]
because its ideal is $(y^4,y^3z,y^2z^2,z^3)$.
A finite module over a discrete valuation ring is
the sum of its free and torsion parts; equality of
generic rank and residue-space dimension here makes
the torsion part zero. Thus $J$ is $x$-saturated.
By (A.T.15), $J$ and $I^2$ agree after inverting $x$.
The curve is lci away from the origin, so its symbolic
square differs from its ordinary square only by a
module supported there. Hence
$I^{(2)}=I^2:x^\infty=J$.

On its normalization let
$a_1=dg_1$ and $b_1=dg_2+t\,dg_1$.
They form a basis of the normalized conormal image;
$dg_3=-2t^2a_1+t b_1$.
The quadratic term of $F$ in this basis is
\[
 q_F=\frac3t a_1^2-\frac3{t^2}a_1b_1
                           +\frac1{t^3}b_1^2,
 \qquad \det(q_F)=\frac3{4t^4}.                  \tag{A.T.16}
\]
Every equation in $I^2$ has a holomorphic quadratic
form in this same basis. If $f_0=cF+I^2$ had
$c(0)\ne0$, the determinant of its quadratic form
would have the noncancellable leading pole
$3c(0)^2/(4t^4)$: the cross terms with a holomorphic
form have pole at most three. It could not have
generic rank one. Thus $c(0)=0$, and (A.T.15) gives
$f_0\in I^2$.

For the second ring put $G=y^2-x^3$ and
$P=(z,G)$, $Q_1=(x,y)$. Then $I=P\cap Q_1$ and
\[
 I^{(2)}=P^2\cap Q_1^2=I^2+(zG).                \tag{A.T.17}
\]
Indeed $P^2=(z^2,zG,G^2)$, while $zG,G^2$ already
belong to $Q_1^2$. Multiplication by $z^2$ is
injective modulo $Q_1^2$, so membership of
$z^2A+zGB+G^2C$ in $Q_1^2$ forces $A\in Q_1^2$.
This gives the asserted generating ideal directly.
Furthermore
$x(zG)=(xz)G$, $y(zG)=(yz)G$, and
$z(zG)=(yz)^2-x(xz)^2$ all lie in $I^2$.
On the cusp branch use $a_1=d(xz)=t^2dz$ and
$b_1=dG=-3t^4dx+2t^3dy$.
The quadratic form of $zG$ is $t^{-2}a_1b_1$,
of determinant $-1/(4t^4)$. The same rank-one
argument forces its coefficient to vanish at the
origin, and the three identities again give
$f_0\in I^2$.

For clarity, the bundle used at these non-lci points
is now defined by its actual morphism. Put
\[
 K=\operatorname{im}\bigl(\nu^*(I/I^2)
                         \longrightarrow\nu^*\Omega_Y^1\bigr),
 \qquad N=K^*.
\]
The coherent image $K$ is locally free of rank two
on the smooth curve $\Sigma$. Let $R_{\rm ram}$ be
the zero divisor of $d\nu$, and let $K_{\rm sat}$
be the kernel of the surjection
$\nu^*\Omega_Y^1\to\Omega_\Sigma^1(-R_{\rm ram})$.
The torsion quotient $K_{\rm sat}/K$ defines a divisor
$J_{\rm jac}$. Taking determinants in these actual
maps gives
\[
 \deg N=-2+\deg(J_{\rm jac}-R_{\rm ram}).         \tag{A.T.18}
\]
For $(t^3,t^4,t^5)$, the wedge $a_1\wedge b_1$
has coefficients $t^8(3,4t,5t^2)$, while $d\nu$
has common factor $t^2$. Thus the correction is
$8-2=6$. For the cusp branch in the second case,
the wedge is $-t^5(2,3t,0)$ and the tangent common
factor is $t$, giving correction four. On its
smooth branch the wedge is $s^2(dx\wedge dy)$,
giving correction two. In both cases (A.T.18)
therefore gives $\deg N=4$.

Because $f_0\in I^2$, its quadratic form is a
holomorphic section of $\operatorname{Sym}^2K\otimes M$,
independent of its expression as a sum of products
of ideal generators. This follows either from its
actual Hessian or from injectivity of
$\operatorname{Sym}^2K\to\nu^*\operatorname{Sym}^2\Omega_Y^1$.
Its saturated radical and quotient again give line
bundles $R,Q$, with $q_0=\deg Q\le0$ and
$\deg R=4-q_0$. The generic quartic invariant
is meromorphic in $L^2$, where
$L=M\otimes Q^{-1}\otimes R^{-2}$; now
\[
                        \deg L^2=-16+2q_0.        \tag{A.T.19}
\]

We bound its poles using actual transverse affine
coordinate maps, with the curve parameter held
fixed during every normal differentiation. For the
space cusp use, for $t\ne0$,
\[
 (u,v,t)\longmapsto
 \left(t^3,\ t^4+\frac{v}{3t^5},\
                   t^5-\frac{u}{t^3}+\frac{2v}{3t^4}\right).
                                                        \tag{A.T.20}
\]
This is locally biholomorphic and its two transverse
derivatives are the dual frame to $a_1,b_1$.
Each coefficient of a derivative has pole at most
five. The first normal derivatives of all three
$g_i$ are holomorphic, since
$dg_1=a_1$, $dg_2=b_1-ta_1$ and
$dg_3=-2t^2a_1+t b_1$.
Second normal derivatives have pole at most ten.
For $f_0=\sum c_{ij}g_i g_j$, the cubic Taylor
coefficients consequently have pole at most ten:
the only surviving terms use a first and a second
derivative of the $g_i$, or a first derivative of
$c_{ij}$ and two first derivatives of the $g_i$.
Every fourth derivative of a holomorphic ambient
function in (A.T.20) has pole at most twenty.

Choose a holomorphic change of the $N$ frame adapted
to its saturated radical and quotient. It does not
increase these bounds. The radical cubic is zero
because the generic transversal is $A_3$.
Writing the relevant Taylor terms as
$a\ell^2+b\ell r^2+er^4$ therefore gives a
holomorphic $a$, a mixed cubic $b$ of pole at most
ten, and $e$ of pole at most twenty. The invariant
$\Delta=ae-b^2/4$ has pole at most twenty.
These are the same intrinsic jets on overlaps:
to weighted order four, where $\ell$ has weight
two and $r$ has weight one, an adapted change is
$\ell'=\alpha\ell+u r^2$, $r'=\beta r$.
Direct substitution multiplies $\Delta$ by
$\alpha^{-2}\beta^{-4}$; changes of the defining
section supply the factor $M^2$. Higher weighted
terms and changes of the curve parameter do not
alter these coefficients. This verifies the
transition law for the line $L^2$ in (A.T.19).

There is only one singular preimage $p$ in this case.
Thus $\Delta$ is a holomorphic section of the square
line $L^2(20p)$, whose degree is $4+2q_0$.
As in (A.T.7), the even pole correction preserves
branch parity, so
\[
 \deg\mathcal B\le4+2q_0,
 \qquad \deg\mathcal B\le-2q_0+1.
\]
The second bound again uses A.5.17 away from that
single singular image preimage. Adding them gives
$2\deg\mathcal B\le5$, contradicting
$\deg\mathcal B=2g+2\ge4$.

For the cusp-plus-line germ, the corresponding
actual affine coordinate maps are
\[
 (u,v,t)\longmapsto
       \left(t^2,t^3+\frac{v}{2t^3},\frac{u}{t^2}\right),
 \qquad
 (u,v,s)\longmapsto\left(\frac us,\frac vs,s\right).
                                                        \tag{A.T.21}
\]
Their first generator derivatives are holomorphic
in the respective normalized conormal frames.
The same product differentiation gives cubic pole
bounds six and two, and the fourth derivative bounds
are twelve and four. Thus $\Delta$ has poles bounded
by $12p+4q$ at the two normalization preimages.
The line $L^2(12p+4q)$ is a square and has degree
$2q_0\le0$ by (A.T.19). A nonzero section has no
positive-degree odd divisor, whereas this double
cover has at least four branch points. This is the
last contradiction. All non-lci cases have now
been excluded, and the positive-genus case at the
start gives the asserted elliptic curves and constant
transversal type.

We finally prove the normalization and Euler assertions,
keeping track of any unnecessary blowups in the original
choice of resolution. Factor $\pi$ through the minimal
resolution $\pi_{\min}:Z_{\min}\to T$, with the resulting
sequence of exceptional-curve blowdowns
$q:Z\to Z_{\min}$. Every curve contracted by $q$ is
vertical for $p$, by the stated assumption on $\pi$.
The equality of structure-sheaf direct images for each
blowdown, in local coordinates on $B$, gives a descended
ruling $p_{\min}:Z_{\min}\to B$ with
$p=p_{\min}q$. The image $C_{\min}$ is still a section,
and the pushed-forward effective conductor is
$D_{\min}=2C_{\min}+V_{\min}$ with $V_{\min}$ vertical.
The normalization-canonical comparison gives
$K_{Z_{\min}}+D_{\min}=\pi_{\min}^*\eta^*\omega_S$
as line bundles, and its rational Chern class is zero.

Statement C2d gives a smooth normalization near the
inverse image of the elliptic curve $D$, with normal
line of degree zero on its unramified double cover.
The minimal resolution is an isomorphism there, so
$C_{\min}^2=0$. No other conductor or exceptional
curve meets it. Hence $V_{\min}\cdot C_{\min}=0$.
No ruling $(-1)$-curve can meet $C_{\min}$: one in
$V_{\min}$ is disjoint from it, while one outside
$V_{\min}$ would have
\[
 (K_{Z_{\min}}+2C_{\min}+V_{\min})\cdot E
       =-1+2C_{\min}\cdot E+V_{\min}\cdot E\ge1,
\]
contrary to the vanishing Chern class.

Contract ruling $(-1)$-curves to a minimal ruled bundle.
Every contraction therefore avoids the section. For
one blowdown $\sigma:W\to W'$ let $D'=\sigma_*D_W$,
let $a$ be the coefficient of its exceptional curve,
and let $d$ be the multiplicity of $D'$ at the centre.
The actual divisor comparison is
\[
 K_W+D_W=\sigma^*(K_{W'}+D')+(1+a-d)E.
                                                        \tag{A.T.22}
\]
Its intersection with $E$ gives $a=d-1$, so the
log-canonical line and its zero rational Chern class
descend. The section remains disjoint from the
effective vertical part at every step. On a minimal
ruled bundle every nonzero effective vertical divisor
is a positive sum of fibres and meets the section.
The final vertical part is consequently zero. If
there had been a blowup in the reverse sequence,
its first centre would be away from this section,
and (A.T.22) would give coefficient $-1$ on its
exceptional curve. Its strict transform retains that
negative coefficient under subsequent blowups, a
contradiction to effectiveness. Thus no ruling
blowdowns were needed: $Z_{\min}$ is already a
minimal ruled bundle and $D_{\min}=2C_{\min}$.

All exceptional curves of $\pi_{\min}$ would be
vertical. A minimal ruled bundle has irreducible
fibres of square zero, which cannot form the
negative-definite exceptional divisor of a surface
resolution. Hence $\pi_{\min}$ has no exceptional
curve and is an isomorphism. This proves the assertion
about $T$. Finally $\eta$ is an isomorphism away from
its conductor section $C_T$ and the elliptic image
$D$. Additivity of the topological Euler number on
these actual decompositions gives
\[
 e(S)=e(T)-e(C_T)+e(D)
      =e(\PP^1)e(B)-0+0=0.                       \tag{A.T.23}
\]
The full multiple fibre has the same underlying
topological space and therefore the same Euler number.
\end{proof}

\paragraph{Connected double covers in the two non-lci cases.}
The non-lci calculation also excludes a connected
degree-two conductor cover when its covering curve
is rational, for either of the two image rings in
(A.T.14). In the cusp-plus-line case the cleared
quartic has degree $2q_0\le0$, so it has no branch
divisor at all. In the space-cusp case a rational
double cover has branch degree two. The already
proved inequalities give
\[
                2\le4+2q_0,\qquad 2\le-2q_0+1,
\]
so the integer $q_0$ is $-1$. The cleared quartic
then has degree two and exactly two simple branch
points. Only one point lies above the singular
image germ. The other is a smooth simple branch
point, excluded by the actual coordinate reduction
(A.T.10)--(A.T.12). Thus these two non-lci images
cannot support a connected double cover of any
nonnegative genus in a torus degeneration.

The integer branch budgets, the discriminant square
weights and the coordinate divisibility used in
(A.T.10)--(A.T.12) are checked in
\path{code/E_tacnode_bounds.py}; the non-lci charts,
symbolic generators and cubic/quartic pole orders
are checked in \path{code/E_low_delta_tacnode.py}.
The topological
quotient (A.T.2) targets the actual singular curve
$D/J$; replacing it by its normalization would
lose the free-group obstruction in (A.T.4).


% A.5.19 from part_A6_triple_images.tex



% Complete checked source block A.5.19; source SHA256 c4d2b5ae417e36d9ce77076d8b673bf64ea5b81792ae8f8a330ab7cfd85a5404
\begin{statement}{A.5.19}{Proved}{Ordinary-triple images have arithmetic genus at most one}
Let $f:Y\to C_0$ be a proper holomorphic map from a smooth
compact complex threefold with $b_2(Y)=0$, with connected
fibres and smooth general fibre a complex two-torus.
Let $mS$ be a full fibre with irreducible reduced support,
and let $D\subset\operatorname{Sing}S$ be an irreducible
curve with generic ordinary-triple plane transversal.
Then $D$ is a local complete intersection in $Y$ and
$p_a(D)\le1$. The possibilities for its underlying curve
are a smooth rational curve, a rational curve with one
ordinary node or cusp, and a smooth elliptic curve.
In the elliptic case the transversal type is constant.

When $D$ is singular rational, every component of the
normalized conductor curve dominating $D$ is rational,
and its map to $D^\nu\simeq\PP^1$ has degree one.
Consequently, in the irrational ruled setup of A.5.18
with the generic transversal replaced by an ordinary
triple point, the image of the noncontracted horizontal
section is either smooth rational or smooth elliptic.
In the elliptic case the horizontal cover is unramified
of degree three, the normalization is smooth minimal
elliptic ruled, and the fibre has Euler number zero.
In the rational case the horizontal cover has degree
two or three; these cases are not excluded by this
statement.
\end{statement}
\begin{proof}
By A.5.14, $p_a(D)\le2$. If its normalization has
positive genus, that statement and A.5.5 already give
a smooth elliptic curve of constant transverse type.
Suppose its normalization $\nu:\Sigma\to D$ is
$\PP^1$, with total delta invariant $\delta\le2$.

First let $D$ be lci. Put
$N=(\nu^*(I_D/I_D^2))^*$ and
$M=\nu^*(\mathcal O_Y(S)|_D)$.
The quadratic normal jet vanishes generically, so
its vanishing descends on the reduced curve through
the locally free graded piece $I_D^2/I_D^3$.
Thus the initial cubic is an actual section of
$\operatorname{Sym}^3N^*\otimes M$.
Adjunction and the curve conductor give
$\deg N=-2+2\delta$, while $\deg M=0$.
Its nonzero binary-cubic discriminant is a section
of $M^4\otimes(\det N)^{-6}$, of degree
\[
                         12-12\delta.             \tag{A.U.1}
\]
Hence $\delta\le1$.

We exclude every non-lci germ of delta at most two
by its symbolic cube. The one- and two-branch analytic
classification is proved in (A.T.14). There is one
additional possibility when three branches occur:
delta at least two is forced by identifying their
constants, so equality makes the ring exactly
\[
 \{(a,b,c)\in\C\{t_1\}\oplus\C\{t_2\}\oplus\C\{t_3\}:
                      a(0)=b(0)=c(0)\}.
\]
Its actual embedding in the smooth threefold is the
union of the three coordinate axes, with ideal
$I=(xy,xz,yz)$. Indeed the three branch parameters
give a basis of the cotangent space, so their lifts
are ambient coordinates as in A.5.18.
Thus there are exactly three non-lci cases to consider;
each consumes the whole delta budget at one point.

For the space cusp retain $I=(g_1,g_2,g_3)$ and
$F$ from (A.T.15). Its exact symbolic-cube identity is
\[
                         I^{(3)}=I^3+FI.           \tag{A.U.2}
\]
Here is a verification in the convergent local ring.
The right-hand side agrees with $I^3$ after inverting
$x$, by (A.T.15). Its quotient is finite over
$\C\{x\}$, of generic rank $3\cdot6=18$: a square
normal plane has six monomials of total degree less
than three. Modulo $x$, its ideal in $\C\{y,z\}$ is
\[
 (y,z)^6+(y^2z^3,yz^4,z^5).
\]
The fifteen monomials of total degree at most four,
together with $y^5,y^4z,y^3z^2$, give residue dimension
eighteen. The same discrete-valuation-module argument
used for (A.T.15) makes the quotient $x$-torsion-free.
It is therefore the saturation of $I^3$, which equals
the symbolic cube since the curve is lci off the origin.
This proves (A.U.2).

For the cusp-plus-line ideal $I=(xz,yz,G)$ with
$G=y^2-x^3$, one has
\[
                   I^{(3)}=I^3+(zG)I.             \tag{A.U.3}
\]
To verify this directly, write $P=(z,G)$ and
$Q=(x,y)$. Their powers are primary because these
prime ideals are complete intersections. The
symbolic cube is $P^3\cap Q^3$.
Write an element of $P^3$ as
$z^3A+z^2GB+zG^2C+G^3E$.
The last two terms are already in $Q^3$.
Reducing the other terms modulo $Q^3$ gives
$zA+GB=0$ there. The constant-in-$(x,y)$ coefficient
of $B$ must therefore be divisible by $z$; write
$B=zB'+B_Q$, $B_Q\in Q$.
It follows that $A+GB'\in Q^3$.
Consequently
\[
 P^3\cap Q^3=(z^3Q^3,z^2GQ,zG^2,G^3)
                        =I^3+(zG)I,
\]
as claimed.

For the three axes, membership of
$x^a y^b z^c$ in the symbolic cube means
$a+b,a+c,b+c\ge3$.
If one exponent is zero, the other two are at
least three and the monomial belongs to $I^3$.
If all are positive, either at least two are at
least two, in which case it belongs to $(xyz)I$,
or two are one, which violates one of the three
inequalities. Hence
\[
                         I^{(3)}=I^3+(xyz)I.       \tag{A.U.4}
\]

In all three cases the generic ordinary-triple
equation $f_0$ belongs to $I^{(3)}$: on the smooth
locus of each curve branch its order in the two
normal variables is at least three. Define the
bundle by its actual conormal image,
\[
 K=\operatorname{im}\bigl(\nu^*(I/I^2)
                           \to\nu^*\Omega_Y^1\bigr),
 \qquad N=K^*.
\]
For the first two cases (A.T.18) gives $\deg N=4$.
For the three axes, on the branch $(t,0,0)$ the
frame is $t\,dy,t\,dz$, with conormal determinant
order two and no tangent ramification. Each of the
three branches has the same correction, for total
six, so again $\deg N=-2+6=4$.

Use the actual affine maps (A.T.20)--(A.T.21).
For the three axes add
\[
              (u,v,t)\longmapsto(t,u/t,v/t)
                                                        \tag{A.U.5}
\]
on the first branch, and the two analogous permuted
maps on the other branches. In these charts the
first normal derivatives of the ideal generators
are holomorphic in the specified conormal frame.
An $I^3$ term consequently has holomorphic initial
cubic. The extra term in each of (A.U.2)--(A.U.4)
has initial cubic equal to the quadratic term of
its displayed symbolic-square generator, multiplied
by the first derivative of an ideal generator.

For the space cusp that quadratic term is (A.T.16),
so the cubic pole is at most three. For the
cusp-plus-line germ it is $t^{-2}uv$ on the cusp
and $s^{-1}v^2$ on the smooth branch, so the pole
bounds are two and one. For the three axes it is
$t^{-1}uv$ in (A.U.5), giving pole bound one on
each branch. Coefficients multiplying these terms
are holomorphic; their derivatives do not enter
the initial cubic because the lower normal jets
of the products vanish.

The initial cubic is intrinsic on the dense smooth
part of $D$, as all lower jets are zero. Its
transition under a normal-coordinate change is
therefore the actual cubic tensor in
$\operatorname{Sym}^3N^*\otimes M$; higher-order
coordinate terms make no contribution. The binary
discriminant is homogeneous of degree four in
these coefficients. It is a meromorphic section
of $M^4\otimes(\det N)^{-6}$, of degree $-24$,
with pole divisor of total degree at most twelve
in each of the three cases. After this exact pole
correction its line has negative degree, at most
$-12$. It cannot have the required nonzero
holomorphic section. This excludes all non-lci
cases and proves that $D$ is lci with $p_a(D)\le1$.
The list of underlying curves follows from the
delta-one classification in A.5.14.

We give the finite-cover comparison at a singular
rational image. Here $\delta=1$, so (A.U.1) has
degree zero and its nonzero section is nowhere zero
on $\Sigma=\PP^1$. The cubic defines a subscheme
$R_3\subset\PP(N)$ of tangent-root lines. Every
fibre consists of three distinct points. In a
projective coordinate chart, its derivative with
respect to the root coordinate is nonzero, so the
implicit function theorem makes
$R_3\to\Sigma$ a finite etale map of degree three.
There is no component in a whole projective fibre,
since its nonzero discriminant forbids an identically
zero cubic there.

Over the generic point of $D$, normalization of an
ordinary triple transversal separates exactly its
three tangent lines. This supplies an isomorphism
between the function-field algebra of $R_3$ and
that of the normalized conductor curves dominating
$D$, after their maps are factored through $\Sigma$.
Both are normal finite curves over $\Sigma$.
Taking the integral closure of $\mathcal O_\Sigma$
in this same algebra therefore identifies the two
finite covers everywhere. Every connected finite
etale cover of the simply connected Riemann sphere
has degree one. Thus each of these conductor
components is rational and maps isomorphically to
$\Sigma$, proving the second assertion.

In the irrational ruled setting the horizontal
section has positive genus. Its finite map to a
rational image normalization must have degree at
least two. The preceding comparison excludes a
singular rational image, and the generic triple
transversal bounds its degree above by three.
The remaining rational image is therefore smooth,
with degree two or three. If the image is elliptic,
no vertical rational curve can dominate it, so all
three branches belong to the horizontal section.
Statement C2d gives an unramified degree-three
cover, normal degree zero and constant type.
The exact minimal-resolution factorization and
blowdown argument (A.T.22) applies without change
to this coefficient-two section: its proof uses
only those properties. It makes the normalization
smooth minimal elliptic ruled with conductor twice
the section. The actual normalization Euler formula
(A.T.23) then gives zero.
\end{proof}

The singular rational possibilities cannot be
removed by a blanket local spherical-cycle argument.
For example
\[
                     f=x(yz)(x-yz)
                                                        \tag{A.U.6}
\]
has ordinary-triple transversals along the nodal curve
$D=(x,yz)=0$ and an additional ordinary-double
critical line $y=z=0$. For $t\ne0$, put
$B_t=\{(x,v):xv(x-v)=t\}$. The exact maps
\[
 B_t\times\C^*\longrightarrow f^{-1}(t),
 \quad (x,v,y)\longmapsto(x,y,v/y),\qquad
 (x,y,z)\longmapsto(x,yz,y)
                                                        \tag{A.U.7}
\]
are inverse holomorphic maps. Nonzero $t$ forces
$v,y$ and $z$ to be nonzero. The smooth affine
curve $B_t$ and $\C^*$ are aspherical, so this
regular affine fibre has no spherical homotopy
class. This is a local model, not a compact
$b_2=0$ example. A deformation smoothing $yz$
changes the additional critical line's transverse
type on the boundary and does not satisfy the
comparison used in A.5.17.

Exact ideal, normal-coordinate and discriminant
calculations are retained in
\path{code/E_triple_images.py} and
\path{code/E_triple_images.out}; the independently
computed symbolic saturations are also in
\path{code/A_low_delta_cusp.py}.


% A.5.20 from part_A6_cusp_completion.tex



% Complete checked source block A.5.20; source SHA256 d3dc59a72d2cd8cc2a598d0b3a0aa60cb8e1826a23973560f266aab63424e1fd
\begin{statement}{A.5.20}{Proved}{A noncontracted horizontal cusp has zero Euler number}
Let $f:Y\to C_0$ be a proper holomorphic map from a smooth
compact complex threefold with $b_2(Y)=0$, with connected
fibres and smooth general fibre a complex two-torus.
Let $mS$ be a full fibre with irreducible reduced support,
$\eta:T\to S$ its normalization and $\pi:Z\to T$ an
irrational ruled resolution with ruling $p:Z\to B$.
Suppose every exceptional curve of $\pi$ is vertical
and the resolution conductor is $2C+V$, where $C$
is a noncontracted section and $V$ is vertical.
If the generic transversal at $D=\eta\pi(C)$ is an
ordinary cusp, then $D$ and $C$ are smooth elliptic
curves and $C\to D$ is an isomorphism. The cusp type
is constant along $D$. The normalization $T$ is a
smooth minimal ruled surface over this elliptic curve,
its conductor is twice a section, and $e(S)=0$.
This permits an initially singular normalization;
smoothness is part of the conclusion.
\end{statement}
\begin{proof}
A transverse cusp has exactly one normalization
branch. The map from the horizontal section to the
normalization of its image consequently has degree
one. Both curves are smooth, so this finite map
identifies $C\simeq B$ with $D^\nu$. Its genus is
positive. Statement A.5.14, including its actual
delta-one lci comparison, and A.5.5 therefore make
$D$ itself smooth elliptic and its cusp transversal
constant at every point. Thus $C\to D$ is an
isomorphism, without any assumption about the
ambient Picard pullback.

For the exact local normalization $x=u^3,y=u^2$,
let $C_T\subset T$ be the image of $C$. The
normalization is smooth near this curve. Its normal
line $P=N_{C_T/T}$ has $P^2=R$ and $P^3=Q$
in the degree-zero radical and Hessian quotient
lines of C2d. Hence $\deg P=0$. The normalization
is smooth near $C$, and no other conductor curve
meets it. Factor $\pi$ through its minimal
resolution, descending the ruling across its
vertical exceptional blowdowns as in A.5.18.
On that minimal resolution the section is isomorphic
to $C_T$ in its unchanged smooth neighbourhood. It
therefore has self-intersection zero and is disjoint from
the effective vertical conductor. The canonical
comparison is the actual line-bundle identity
$K+2C+V=(\eta\pi)^*\omega_S$, of zero rational
Chern class.

Apply (A.T.22) to each ruling blowdown. Its
exceptional coefficient is exactly the multiplicity
at its centre minus one. No ruling $(-1)$-curve
can meet $C$, since the intersection with $K+2C+V$
would then be positive. On the minimal ruled
bundle the vertical conductor is zero because it
is disjoint from a section. A first reverse blowup
away from that section would have coefficient
$-1$, so there were no ruling blowups. All
resolution-exceptional curves are vertical, and
the square-zero fibres of the resulting bundle
cannot be exceptional. This identifies the minimal
resolution with $T$ itself and proves the asserted
smoothness and minimality.

Finally the normalization is an isomorphism off
$C$ and $D$. The exact Euler comparison is
$e(S)=e(T)-e(C)+e(D)$; all three terms vanish,
since $T$ is ruled over an elliptic curve and
$C,D$ are elliptic. Thus $e(S)=0$, also for the
underlying topological space of the full fibre $mS$.
\end{proof}


% A.5.21 from part_A6_binary_products.tex



% Complete checked source block A.5.21; source SHA256 9ff0ac5af26dc04cfc558d7ddbead982a51bc89e1074cf6ec4f2e770c96ed9f4
\begin{statement}{A.5.21}{Proved}{A homogeneous binary-product grid is excluded}
Let $f:Y\to C_0$ be a proper holomorphic map from a smooth
complex threefold with smooth nearby fibres complex
two-tori. Let $P(X,V)$ be a homogeneous polynomial of
degree $d\ge2$, a product of distinct complex linear
forms, and suppose $P(1,0)\ne0$.
There is no point at which the actual fibre map has
the local equation
\[
             f=u(x,y,z)\,P(x,yz)^m,
             \qquad m\ge1,                         \tag{A.BP.1}
\]
with $u$ a holomorphic unit. In particular this excludes
$\prod_{i=1}^d(x-a_i yz)$ for distinct constants $a_i$,
and every positive full-fibre multiplicity of that germ.
No Betti-number assumption on $Y$ is required.

The same exclusion holds for the variable-coefficient
equation $f=u h^m$ with
\[
 h=\sum_{i=0}^d a_i(x,y,z)x^{d-i}(yz)^i,
 \qquad P_0(X,V)=\sum_{i=0}^d a_i(0)X^{d-i}V^i,
                                                        \tag{A.BP.4}
\]
provided $P_0$ has $d$ distinct projective roots and
$P_0(1,0)\ne0$. The comparison with its constant
initial form is the explicit analytic family below.
\end{statement}
\begin{proof}
Put $h=P(x,yz)$. It is weighted homogeneous of degree
$2d$ for the positive weights $(2,1,1)$.
On a small simply connected coordinate neighbourhood
take a holomorphic logarithm of the unit and set
\[
 \Phi(x,y,z)=\bigl(xu^{1/(md)},
                  yu^{1/(2md)},zu^{1/(2md)}\bigr).
                                                        \tag{A.BP.2}
\]
Its derivative at the origin is invertible, and direct
homogeneity gives $h\Phi=u^{1/m}h$. Thus
$f=(h\Phi)^m$ exactly. It suffices to construct a
compact spherical cycle of nonzero square in $h=1$,
because weighted homotheties then put it inside an
arbitrarily small actual fibre of (A.BP.1).

The curve $B=\{P(x,v)=1\}\subset\C^2$ is smooth:
Euler's identity $xP_x+vP_v=dP$ precludes a common
zero of its two partial derivatives on that level.
It is also connected. The map
\[
 B\longrightarrow\PP^1\setminus\{P=0\},
                       \quad(x,v)\longmapsto[x:v]
\]
is an etale cover of degree $d$. Locally a projective
representative has nonzero $P$-value, and its $d$
rescalings with $P=1$ give the sheets. Around any
one of the simple projective zeros of $P$, that
value has a simple zero, so its inverse $d$th root
has monodromy a full $d$-cycle. This monodromy is
transitive on the sheets, proving connectedness.

The function $v$ on $B$ has exactly $d$ distinct
simple zeros, since
$P(x,0)=P(1,0)x^d$ and $P_x(x,0)=d/x$ at those
points. Choose an embedded arc $\gamma$ in $B$
joining two of its zeros and avoiding all other
zeros. On the interior choose a continuous argument
of $v(\gamma(t))$ and two functions $A(t),B_0(t)$
with $A B_0=v(\gamma(t))$ and
$|A|=|B_0|=|v(\gamma(t))|^{1/2}$.
The map
\[
 (t,\theta)\longmapsto
   \bigl(x(\gamma(t)),A(t)e^{\mathrm i\theta},
                        B_0(t)e^{-\mathrm i\theta}\bigr)
                                                        \tag{A.BP.3}
\]
lies in $h=1$ and collapses both boundary circles
of $[0,1]\times S^1$ to points. It is therefore a
spherical cycle, with embedded interior circles.

A parallel-arc isotopy with the same endpoints gives
a homologous spherical cycle whose support meets
the first only at those endpoints. At each endpoint,
$P_x\ne0$, so $(y,z)$ are actual holomorphic
coordinates on $h=1$. The links of the two cycles
are parallel curves of slope $(1,-1)$ on a Clifford
torus in the local three-sphere. Their linking
number is $-1$. As in the explicit relative-chain
comparison in A.5.15, that number is the local
intersection of the two coned cycles. Both
orientations reverse at the second endpoint,
so its intersection sign is unchanged. The
spherical class consequently has square $-2$.

The weighted homothety
$(x,y,z)\mapsto(\lambda^2x,\lambda y,\lambda z)$
maps $h=1$ biholomorphically to $h=\lambda^{2d}$
and shrinks the compact cycles and their isotopy
to the origin. Together with (A.BP.2), this puts
them inside an actual smooth global fibre of $f$.
That fibre is a complex two-torus and has zero
$\pi_2$, so the spherical class has zero homology
image. The oriented codimension-zero embedding
preserves the intersection of the compact cycles,
giving both zero and $-2$ for its square. This
contradiction proves the statement.

We now prove the variable-coefficient assertion,
including the boundary comparison. Absorb a local
$m$th root of the unit into $h$; this preserves the
form (A.BP.4) and its two initial-form conditions.
For a complex parameter $s$ set
\[
 h_s=\sum_{i=0}^d a_i(s^2x,sy,sz)x^{d-i}(yz)^i.
                                                        \tag{A.BP.5}
\]
This is holomorphic also at $s=0$, with
$h_0=P_0(x,yz)$. For $s\ne0$ there is the exact
identity
\[
       h(s^2x,sy,sz)=s^{2d}h_s(x,y,z).           \tag{A.BP.6}
\]
The leading coefficient $a_0$ is a unit. The implicit
function theorem applied to its degree-$d$ root
polynomial gives distinct holomorphic functions
$b_i(s,x,y,z)$ and the exact factorization
\[
       h_s=a_0(s^2x,sy,sz)
                     \prod_{i=1}^d(x-b_i(s,x,y,z)yz)
                                                        \tag{A.BP.7}
\]
on one small parameter and coordinate neighbourhood.
Each factor defines a smooth surface, since its
$x$ derivative is a unit. For $i\ne j$ the ideal
generated by the two factors equals $(x,yz)$:
their difference is the unit $b_j-b_i$ times $yz$.
Thus the zero level has precisely the same two
critical curves $D=(x,yz)=0$ for every $s$.

Choose a small closed ball $K$ in the coordinate
chart such that the two axes of $D$ meet its boundary
transversely and the zero-level surfaces away from
those axes are transverse to the boundary. At
$s=0$ this follows directly for the smooth graph
surfaces $x=b_i(0)yz$ and sufficiently small radii.
It persists for small $s$ by compactness away from
neighbourhoods of the two boundary links.
Near either link, use normal coordinates $(U,V)$
for the smooth critical curve and a parameter $w$
along it. The normal initial term of $h_s$ is a
degree-$d$ binary form with distinct roots.
On the compact parameterized link its normal
gradient therefore satisfies a uniform bound
\[
 \|\partial_{U,V}h_s\|\ge c\|(U,V)\|^{d-1},
 \qquad |\partial_w h_s|\le C\|(U,V)\|^d
                                                        \tag{A.BP.8}
\]
after the neighbourhood is shrunk. The first bound
follows because the gradient of a homogeneous
reduced binary form has no zero on the unit sphere;
higher normal terms have one additional order.
The second follows from $h_s\in(U,V)^d$.

Transversality of the critical curve to $\partial K$
allows a boundary coordinate with one real component
of $w$ solved as a smooth function of $(U,V)$ and
the remaining component. The correction to normal
differentiation tangent to this boundary is a
bounded multiple of $\partial_w h_s$.
By (A.BP.8) it is smaller by one normal order than
the surjective complex normal derivative. Thus
$h_s|_{\partial K}$ is a real submersion at every
nearby point with $h_s\ne0$, uniformly in $s$.
Away from the two link neighbourhoods, compactness
and the preceding zero-level transversality give
the same conclusion for all sufficiently small
nonzero values.

Remove the proper critical-value discriminant from
the product of the parameter disc and a sufficiently
small punctured value disc. The critical locus over
that domain cannot meet the boundary, so its
projection is proper and its image is complex
analytic. Its complement is connected. The level
sets in $K$ form a proper submersion with boundary
there, and Ehresmann transport identifies their
compact fibres by oriented diffeomorphisms.
Choose a small sphere cycle of the already proved
model $h_0=P_0(x,yz)$, and transport it to a fibre
of $h_{s_0}$ for $s_0\ne0$. Equation (A.BP.6)
then embeds that fibre into the actual fibre
of $h^m$ with value $s_0^{2dm}\epsilon^m$.
It carries the spherical class of square $-2$
into a global torus. The same intersection and
asphericity contradiction proves (A.BP.4).
\end{proof}

The condition $P(1,0)\ne0$ is used to produce the
distinct collapsing circles. The exact product
model (A.U.6), for which $P(X,V)=XV(X-V)$, has
$P(1,0)=0$ and is outside this statement.
The variable-coefficient extension uses only the
specific comparison (A.BP.5)--(A.BP.8), with its
unchanged boundary critical curves. The weighted scaling,
simple-zero and endpoint-square identities are
checked in \path{code/E_binary_products.py}.

\paragraph{Two exact scalar-zero comparisons.}
The homogeneous germ
\[
                  h=xy\{z(x-y)+xy\}
\tag{A.BP.9}
\]
has a cubic normal tensor along $(x,y)=0$ which
vanishes identically at $z=0$. On $xy\ne0$ put
\[
 v=x-y,\qquad Z=y(y+v)z+y^2(2y+v).
\]
These are actual coordinates with $y$, with inverse
$x=y+v$, $z=(Z-y^2(2y+v))/(y(y+v))$, and
\[
                              h=vZ+y^4.
\tag{A.BP.10}
\]
Choose an arc between two roots of $y^4=1$ avoiding
$y=0$ and the other roots. The matching circles
$v=\rho A e^{\mathrm i\theta}$,
$Z=\rho^{-1}B e^{-\mathrm i\theta}$, where
$AB=1-y^4$ and $|A|=|B|=|1-y^4|^{1/2}$,
collapse at its endpoints. Choose the positive
constant $\rho$ small enough that $|v|<|y|/2$
on the arc. The entire spherical cycle then lies
in the indicated coordinate chart. Its two endpoint
linkings are $-1$, as in (A.BP.3), so its square
is $-2$. Ordinary homothety in $(x,y,z)$ shrinks
it because (A.BP.9) has degree four. The same
weighted unit map as (A.BP.2), with all three
weights equal to one, includes every local unit
and every positive full-fibre multiplicity.
Thus this particular scalar-zero germ is excluded.

The following different scalar-zero model records
the limit of that argument:
\[
                     h=xz(y^2-x^2z).
\tag{A.BP.11}
\]
For $h=t\ne0$, the substitution $u=y/x$ is
invertible and gives $h=x^3 z(u^2-z)$.
The map forgetting $x$ is a finite etale triple
cover of $E=\{z(u^2-z)\ne0\}$.
There is an explicit finite etale double cover
\[
 \C^*\times(\C\setminus\{1,-1\})\longrightarrow E,
                  \qquad(s,v)\longmapsto(s^2,sv)
\tag{A.BP.12}
\]
in coordinates $(z,u)$, with deck involution
$(s,v)\mapsto(-s,-v)$. Its source is aspherical,
so both $E$ and the regular affine fibre of
(A.BP.11) are aspherical by these actual covering
maps. The quadratic surface branch has normalization
$(v,u)\mapsto(v,vu,u^2)$ and conductor
$2(v=0)+2(u=0)$ in the full reduced surface.
The other two normalization branches have conductors
$2(y=0)+(z=0)$ on $x=0$, and
$(x=0)+2(y=0)$ on $z=0$.
Thus their reduced conductor supports have only
normal-crossing pairs. This is an exact local
model, not a construction of a compact $b_2=0$
fibre. It shows why a vanishing cubic coefficient
alone does not justify a local spherical-cycle
exclusion. The identities and the two finite-cover
formulas are checked in \path{code/E_cubic_double.py}.

